%% dissertacao_tomo_sar.tex
%% Dissertação de Doutorado
%% Título: Otimização da Trajetória do Receptor em SAR Biestático
%%         com Voo Helicoidal — Exploração do Ângulo de Brewster
%%         com Restrição de Resolução Espacial
%% Idioma: Português
%% Compilar: pdflatex -interaction=nonstopmode dissertacao_tomo_sar.tex  (duas passagens)

\documentclass[12pt,a4paper,oneside]{book}

%% ─── Pacotes ─────────────────────────────────────────────────────────────────
\usepackage[utf8]{inputenc}
\usepackage[T1]{fontenc}
\usepackage[portuguese]{babel}
\usepackage{lmodern}
\usepackage{geometry}
\geometry{left=3cm, right=2.5cm, top=3cm, bottom=2.5cm}

\usepackage{amsmath,amssymb,amsfonts,amsthm}
\usepackage{mathtools}
\usepackage{bm}
\usepackage{siunitx}
\sisetup{output-decimal-marker={,}}

\usepackage{graphicx}
\graphicspath{{figures/}}

\usepackage{float}
\usepackage{caption}
\usepackage{subcaption}
\usepackage[table,xcdraw]{xcolor}
\usepackage{booktabs}
\usepackage{array}
\usepackage{multirow}
\usepackage{longtable}
\usepackage{setspace}
\onehalfspacing

\usepackage{tcolorbox}
\tcbuselibrary{theorems,skins,breakable}

\usepackage[unicode,colorlinks=true,
  linkcolor=blue!60!black,
  citecolor=green!50!black,
  urlcolor=cyan!60!black,
  pdftitle={Otimizacao da Trajetoria do Receptor em SAR Bistatico Helicoidal},
  pdfauthor={David Atencia}]{hyperref}

\usepackage{fancyhdr}
\setlength{\headheight}{15pt}
\pagestyle{fancy}
\fancyhf{}
\fancyhead[L]{\small\nouppercase{\leftmark}}
\fancyhead[R]{\small\thepage}
\renewcommand{\headrulewidth}{0.4pt}

%% ─── Teoremas adicionais ─────────────────────────────────────────────────────
\theoremstyle{remark}
\newtheorem*{observacao}{Observação}

%% ─── Operadores ──────────────────────────────────────────────────────────────
\DeclareMathOperator{\sinc}{sinc}
\newcommand{\vect}[1]{\boldsymbol{#1}}
\newcommand{\uvec}[1]{\hat{\boldsymbol{#1}}}
\newcommand{\dd}{\mathrm{d}}
\newcommand{\jj}{\mathrm{j}}
\newcommand{\Real}{\mathbb{R}}

%% ─── Teoremas ────────────────────────────────────────────────────────────────
\newtheorem{teorema}{Teorema}[chapter]
\newtheorem{proposicao}{Proposição}[chapter]
\newtheorem{lema}{Lema}[chapter]
\newtheorem{corolario}{Corolário}[chapter]
\theoremstyle{definition}
\newtheorem{definicaom}{Definição}[chapter]

%%═══════════════════════════════════════════════════════════════════════════════
\begin{document}
%%═══════════════════════════════════════════════════════════════════════════════

%%─── Capa ─────────────────────────────────────────────────────────────────────
\frontmatter
\begin{titlepage}
  \centering\vspace*{1.5cm}
  {\Large\textsc{Universidade Estadual de Campinas}\par}
  {\large\textsc{Faculdade de Engenharia Elétrica e de Computação}\par}
  \vspace{0.5cm}
  {\large\textsc{Programa de Pós-Graduação em Engenharia Elétrica}\par}
  \vfill
  {\Huge\bfseries Otimização da Trajetória do Receptor\\[0.4em]
   em Radar de Abertura Sintética Biestático\\[0.4em]
   com Voo Helicoidal\par}
  \vspace{1cm}
  {\Large\itshape Exploração do Ângulo de Brewster com\\
   Restrição de Resolução Espacial\par}
  \vfill
  {\large\textbf{David Atencia}\par}
  \vspace{0.6cm}
  {\normalsize Tese apresentada à Faculdade de Engenharia Elétrica e de
   Computação da Universidade Estadual de Campinas como parte dos requisitos
   exigidos para a obtenção do título de Doutor em Engenharia Elétrica.\par}
  \vspace{0.4cm}
  {\normalsize Campinas, 2026\par}
\end{titlepage}

%%─── Dedicatória ──────────────────────────────────────────────────────────────
\cleardoublepage
\thispagestyle{empty}
\vspace*{\fill}
\begin{flushright}
  % TODO: Escrever dedicatória.
  \textit{[Dedicatória a ser escrita.]}
\end{flushright}
\vspace*{\fill}

%%─── Agradecimentos ───────────────────────────────────────────────────────────
\chapter*{Agradecimentos}
\addcontentsline{toc}{chapter}{Agradecimentos}

% TODO: Escrever agradecimentos (orientador, banca, agências de fomento,
%       colaboradores, família).

\textit{[Agradecimentos a serem escritos.]}

%%─── Resumo ───────────────────────────────────────────────────────────────────
\chapter*{Resumo}
\addcontentsline{toc}{chapter}{Resumo}

Esta tese propõe e valida um método de otimização da trajetória do receptor em
um sistema de Radar de Abertura Sintética (SAR) biestático com geometria de voo
helicoidal cônico, orientado ao sensoriamento tridimensional de alvos
subsuperficiais.

O problema é formulado como uma otimização multi-objetivo que maximiza
simultaneamente a potência recebida e a resolução espacial tridimensional.
A solução explora o \emph{ângulo de Brewster} da polarização TM (modo
transverso-magnético), fenômeno puramente eletromagnético pelo qual a
reflectância na interface ar-solo se anula, transmitindo toda a energia
incidente para o subsolo. A trajetória ótima do receptor é aquela cujo ângulo
de incidência na interface é, para cada posição do transmissor, o mais próximo
possível do ângulo de Brewster, sujeito às restrições de resolução.

O fundamento matemático repousa em quatro pilares: (i) as equações de Maxwell
e as condições de contorno que geram os coeficientes de Fresnel em polarização
TM; (ii) o princípio variacional de Fermat para a determinação do raio refratado
na interface; (iii) a análise do espaço de números de onda (\textit{k}-espaço)
do SAR biestático helicoidal, da qual emergem as fórmulas analíticas de
resolução tridimensional; e (iv) o formalismo de matrizes de transferência ABCD
para a extensão a $N$ camadas subsuperficiais.

A validação por simulação de \textit{backprojection} biestático confirma as
fórmulas analíticas com erros inferiores a 0{,}1\% na resolução vertical
$\delta_z$ e inferiores a 12\% na resolução horizontal $\delta_{xy}$.

\vspace{1em}
\noindent\textbf{Palavras-chave:} SAR biestático; tomografia SAR; ângulo de
Brewster; coeficientes de Fresnel; resolução espacial; $k$-espaço; meios
estratificados; formalismo ABCD.

%%─── Abstract ─────────────────────────────────────────────────────────────────
\chapter*{Abstract}
\addcontentsline{toc}{chapter}{Abstract}

This thesis proposes and validates a method for optimizing the receiver
trajectory in a bistatic Synthetic Aperture Radar (SAR) system with conical
helical flight geometry, aimed at three-dimensional subsurface target sensing.

The problem is formulated as a multi-objective optimization that simultaneously
maximizes received power and three-dimensional spatial resolution. The solution
exploits the \emph{Brewster angle} of TM polarization, a purely electromagnetic
phenomenon by which the reflectance at the air-soil interface vanishes,
transmitting all incident energy into the subsurface. The optimal receiver
trajectory is one whose angle of incidence at the interface is, for each
transmitter position along its helix, as close as possible to the Brewster
angle, subject to resolution constraints.

The mathematical foundation rests on four pillars: (i) Maxwell's equations and
boundary conditions yielding the TM Fresnel coefficients; (ii) Fermat's
variational principle for refracted ray determination; (iii) wavenumber
($k$-space) analysis of the bistatic helical SAR, from which analytical
three-dimensional resolution formulas emerge; and (iv) the ABCD transfer matrix
formalism for extension to $N$ subsurface layers.

Validation via bistatic backprojection simulation confirms the analytical
formulas with errors below 0.1\% in vertical resolution $\delta_z$ and below
12\% in horizontal resolution $\delta_{xy}$.

\vspace{1em}
\noindent\textbf{Keywords:} bistatic SAR; SAR tomography; Brewster angle;
Fresnel coefficients; spatial resolution; $k$-space; stratified media; ABCD
formalism.

%%─── Listas ───────────────────────────────────────────────────────────────────
\tableofcontents
\listoffigures
\listoftables

%%─── Lista de Acrônimos ───────────────────────────────────────────────────────
\chapter*{Lista de Acrônimos}
\addcontentsline{toc}{chapter}{Lista de Acrônimos}

% TODO: Completar com todos os acrônimos do texto.

\begin{longtable}{ll}
\toprule
Acrônimo & Descrição \\
\midrule
ABCD   & Matriz de parâmetros de transmissão (Transferência) \\
BP     & \textit{Backprojection} (retroprojeção) \\
FFBP   & \textit{Fast Factorized Backprojection} \\
FMCW   & \textit{Frequency-Modulated Continuous Wave} \\
GPR    & \textit{Ground Penetrating Radar} \\
InSAR  & \textit{Interferometric SAR} (SAR interferométrico) \\
LFM    & \textit{Linear Frequency-Modulated} (chirp linear) \\
PCHIP  & \textit{Piecewise Cubic Hermite Interpolating Polynomial} \\
PSF    & \textit{Point Spread Function} (função de espalhamento de ponto) \\
RCS    & \textit{Radar Cross Section} (seção transversal de radar) \\
SAR    & \textit{Synthetic Aperture Radar} (radar de abertura sintética) \\
SNR    & \textit{Signal-to-Noise Ratio} (relação sinal-ruído) \\
TE     & Transverso-Elétrico (polarização TE) \\
TIR    & Reflexão Interna Total (\textit{Total Internal Reflection}) \\
TM     & Transverso-Magnético (polarização TM) \\
UWB    & \textit{Ultra-Wideband} (ultra-banda larga) \\
\bottomrule
\end{longtable}

%%─── Lista de Símbolos ────────────────────────────────────────────────────────
\chapter*{Lista de Símbolos}
\addcontentsline{toc}{chapter}{Lista de Símbolos}

\begin{longtable}{lll}
\toprule
Símbolo & Descrição & Unidade \\
\midrule
$c$                                       & Velocidade da luz no vácuo                         & m/s       \\
$\omega = 2\pi f$                         & Frequência angular                                  & rad/s     \\
$k = \omega n/c$                          & Número de onda                                      & rad/m     \\
$\lambda_0 = c/f_0$                       & Comprimento de onda no vácuo                        & m         \\
$B$                                       & Largura de banda chirp                              & Hz        \\
$\varepsilon_r = n^2$                     & Permissividade relativa                             & --        \\
$\eta_i = \eta_0/n_i$                     & Impedância intrínseca do meio $i$                   & $\Omega$  \\
$\theta_i$                                & Ângulo de incidência / refração na camada $i$       & rad       \\
$\theta_B = \arctan(n_2/n_1)$             & Ângulo de Brewster                                  & rad       \\
$r_{\text{TM}}$                           & Coeficiente de reflexão TM (amplitude)              & --        \\
$t_{\text{TM}}$                           & Coeficiente de transmissão TM (amplitude)           & --        \\
$R_{\text{TM}} = |r_{\text{TM}}|^2$       & Reflectância TM (potência)                          & --        \\
$T_{\text{TM}} = 1 - R_{\text{TM}}$       & Transmitância TM (potência)                         & --        \\
$\mathcal{L}$                             & Caminho óptico                                      & m         \\
$p = n_0\sin\theta_0$                     & Invariante de Snell (parâmetro de raio)             & --        \\
$\mathbf{M}_i$                            & Matriz ABCD da camada $i$                           & --        \\
$k_{z,i}$                                 & Número de onda vertical na camada $i$               & rad/m     \\
$Z_i = n_i^2/\sqrt{n_i^2-p^2}$            & Impedância TM da camada $i$                         & --        \\
$\Gamma_{\text{TM}}$                      & Coeficiente de reflexão global (ABCD)               & --        \\
$P_t,\, P_r$                              & Potência transmitida / recebida                     & W         \\
$G_t,\, G_r$                              & Ganhos de antena Tx e Rx                            & --        \\
$\sigma$                                  & Seção transversal de radar (RCS)                    & m$^2$     \\
$R_T,\, R_R$                              & Comprimentos geométricos dos raios refratados       & m         \\
$\psi_0 = \arctan(\rho_0/z_0)$            & Ângulo de \textit{look} médio                       & rad       \\
$\theta_{t,0}$                            & Ângulo de transmissão no solo no ponto médio        & rad       \\
$B_{helix}$                               & Comprimento de arco da hélice                       & m         \\
$\beta = \arctan(\Delta z/\Delta\rho)$    & Ângulo de inclinação da hélice                      & rad       \\
$B_\perp = B_{helix}|\cos(\beta-\psi_0)|$ & Apertura tomográfica efetiva                        & m         \\
$\Delta\phi$                              & NorthOffset Tx--Rx                                  & rad       \\
$W_z$                                     & Extensão espectral vertical (frequência)            & Hz        \\
$\delta_z,\, \delta_{xy}$                 & Resoluções vertical e horizontal                    & m         \\
$J_0$                                     & Função de Bessel de $1^\text{a}$ espécie, ordem 0  & --        \\
$J(\vect{r}_{Rx};\alpha)$                 & Função de custo multi-objetivo                      & --        \\
$\alpha$                                  & Parâmetro de compromisso potência--resolução        & --        \\
\bottomrule
\end{longtable}

%%═══════════════════════════════════════════════════════════════════════════════
\mainmatter
%%═══════════════════════════════════════════════════════════════════════════════

%%═══════════════════════════════════════════════════════════════════════════════
\chapter{Introdução}
\label{cap:intro}
%%═══════════════════════════════════════════════════════════════════════════════

% TODO: Revisar e expandir toda a introdução após completar os demais capítulos.

\section{Contexto e Motivação}
\label{sec:intro_contexto}

O Radar de Abertura Sintética (SAR) é uma modalidade de radar coerente que
processa o eco de múltiplas posições sequenciais da antena para sintetizar uma
abertura virtual de dimensão muito superior à antena física, produzindo imagens
de alta resolução independentes das condições atmosféricas e da iluminação
solar~[1]. Na configuração \emph{biestática}, transmissor (Tx)
e receptor (Rx) operam em plataformas separadas, conferindo ao sistema maior
flexibilidade geométrica e a possibilidade de explorar ângulos de iluminação e
recepção complementares.

Quando ambas as plataformas descrevem trajetórias helicoidais cônicas ao redor
da cena, o sistema realiza tomografia SAR tridimensional completa: a hélice
fornece cobertura tanto na frequência espacial vertical ($k_z$, via apertura
tomográfica) quanto na frequência espacial horizontal ($k_{xy}$, via cobertura
azimutal circular). Esta geometria é particularmente vantajosa para o
sensoriamento de alvos \emph{subsuperficiais} — túneis, cavidades, minas
antipessoais, tubulações, camadas estratigráficas — nos quais a onda
eletromagnética precisa atravessar a interface ar-solo e propagar-se em meio
dielétrico antes de atingir o alvo.

O desempenho do sistema em regime subsuperficial é determinado por duas
grandezas fundamentais e conflitantes entre si: (i) a \emph{potência recebida}
$P_r$, que deve ser suficiente para garantir relação sinal-ruído (SNR) adequada
ao processamento da imagem; e (ii) a \emph{resolução espacial tridimensional}
$(\delta_{xy},\delta_z)$, que determina a capacidade de distinção entre alvos
próximos. A potência recebida depende criticamente do ângulo de incidência na
interface ar-solo, pois os coeficientes de reflexão e transmissão de Fresnel
variam fortemente com esse ângulo. Em polarização TM (transverso-magnética),
existe um ângulo privilegiado — o \emph{ângulo de Brewster} $\theta_B$ — no
qual a reflectância é rigorosamente nula e toda a energia incidente é
transmitida para o subsolo~[5]. Por outro lado, a resolução
horizontal $\delta_{xy}$ depende do ângulo azimutal entre Tx e Rx (NorthOffset
$\Delta\phi$), e a resolução vertical $\delta_z$ depende da largura de banda e
da apertura tomográfica da hélice~[3].

A questão central desta tese é, portanto: \emph{qual trajetória do receptor, para
cada posição do transmissor ao longo de sua hélice, maximiza simultaneamente a
potência recebida e a resolução espacial tridimensional da imagem SAR do alvo
subsuperficial?}

\section{O Sistema de Sensoriamento Proposto}
\label{sec:intro_sistema}

O sistema proposto consiste em dois veículos aéreos não tripulados (VANTs) que
descrevem trajetórias helicoidais cônicas ao redor de uma cena de interesse. O
transmissor segue uma hélice fixa e pré-determinada. O receptor, por sua vez,
segue uma trajetória otimizada que, para cada posição do Tx, posiciona o Rx de
modo a maximizar a função de custo multi-objetivo que combina potência recebida
e resolução.

O alvo encontra-se abaixo da superfície do solo. A interface ar-solo é modelada
como plana e horizontal. O solo é tratado como dielétrico homogêneo sem perdas
na aproximação de referência (Capítulos~\ref{cap:fundamentos}--\ref{cap:resolucao}),
e estendido a $N$ camadas estratificadas no modelo final
(Capítulo~\ref{cap:multicamadas}). A polarização da onda transmitida é TM, de
modo a explorar o ângulo de Brewster da interface ar-solo.

Os blocos funcionais do sistema são:
\begin{itemize}
  \item \textbf{Geometria da hélice cônica}: parametrização de Tx e Rx por raio,
  altura e número de voltas, com NorthOffset $\Delta\phi$ variável.
  \item \textbf{Modelo eletromagnético}: coeficientes de Fresnel TM e ângulo de
  Brewster para a interface ar-solo (uma camada) e formalismo ABCD para $N$
  camadas.
  \item \textbf{Traçado de raios}: princípio de Fermat com solução iterativa de
  Newton-Raphson para determinar o ponto de refração e o comprimento do raio
  refratado para cada par (posição do sensor, alvo).
  \item \textbf{Análise de resolução}: derivação analítica de $\delta_z$ e
  $\delta_{xy}$ a partir do $k$-espaço biestático.
  \item \textbf{Otimizador multi-objetivo}: minimização da função de custo
  logarítmica $J(\vect{r}_{Rx};\alpha)$ para cada posição discreta do Tx,
  seguida de interpolação PCHIP da trajetória completa do Rx.
  \item \textbf{Validação}: simulador de \textit{backprojection} biestático com
  refração para verificação das fórmulas analíticas.
\end{itemize}

\section{Objetivos e Contribuições}
\label{sec:intro_contrib}

O objetivo geral desta tese é desenvolver e validar um método sistemático de
otimização da trajetória do receptor para SAR biestático helicoidal orientado
ao sensoriamento subsuperficial em polarização TM.

Os objetivos específicos, que correspondem às contribuições originais deste
trabalho, são:

\begin{enumerate}

  \item \textbf{Derivação eletromagnética rigorosa dos coeficientes de Fresnel TM
  e do ângulo de Brewster.} A partir das equações de Maxwell e das condições de
  contorno, demonstra-se que $\theta_B = \arctan(n_2/n_1)$ é o único zero da
  reflectância TM para meios sem perdas e não-magnéticos, com interpretação
  física em termos da radiação dipolar.

  \item \textbf{Derivação analítica das equações de resolução 3D do SAR biestático
  helicoidal.} As fórmulas para $\delta_z$ e $\delta_{xy}$ são obtidas pela
  análise do vetor de número de onda biestático refratado e do perfil de Bessel
  $J_0$ da PSF circular, generalizando resultados da literatura para incluir
  o efeito de refração na interface ar-solo.

  \item \textbf{Formulação de uma função de custo multi-objetivo logarítmica para
  a otimização conjunta de potência e resolução.} A escala logarítmica resolve
  o problema de comensurabilidade entre potência ($\sim 10^{-15}$\,W) e resolução
  ($\sim 10^{-4}$\,m$^2$), e a fronteira de Pareto identifica a região de operação
  recomendada.

  \item \textbf{Extensão do modelo de transmissão TM para $N$ camadas
  subsuperficiais via formalismo de matrizes de transferência ABCD.} O modelo
  trata a última camada como semi-infinita (sem reflexão na face inferior),
  com impedâncias de entrada e saída assimétricas ($Z_{\text{in}} \neq Z_{\text{out}}$),
  e reduz-se exatamente ao caso de Fresnel quando $N=1$.

  \item \textbf{Validação por simulação de \textit{backprojection} biestático.}
  O simulador integra numericamente a lei de Snell via Newton-Raphson e confirma
  as fórmulas analíticas com erros inferiores a 0{,}1\% em $\delta_z$ e
  inferiores a 12\% em $\delta_{xy}$.

\end{enumerate}

\section{Organização da Tese}
\label{sec:intro_org}

Esta tese está organizada em sete capítulos e dois apêndices, seguindo uma
estrutura cumulativa em que cada capítulo estabelece os fundamentos necessários
para o seguinte.

O \textbf{Capítulo~\ref{cap:fundamentos}} (Fundamentos Teóricos) consolida a
base teórica do trabalho: as equações de Maxwell e os coeficientes de Fresnel TM,
o ângulo de Brewster, o princípio de Fermat para o traçado de raios refratados,
a geometria da hélice cônica e a equação do radar biestático.

O \textbf{Capítulo~\ref{cap:resolucao}} (Resolução Espacial 3D) deriva as
fórmulas analíticas de resolução $\delta_z$ e $\delta_{xy}$ a partir da análise
do $k$-espaço biestático refratado, e discute a condição ótima de inclinação da
hélice e a correção \textit{off-axis} de Góes~[3].

O \textbf{Capítulo~\ref{cap:multicamadas}} (Meios Estratificados) estende o
modelo eletromagnético do Capítulo~\ref{cap:fundamentos} para $N$ camadas via
formalismo ABCD, demonstra a consistência com o caso de Fresnel e discute o
ângulo de transmissão ótimo numérico para pilhas multicamada.

O \textbf{Capítulo~\ref{cap:otimizacao}} (Otimização da Trajetória) formula o
problema de otimização multi-objetivo sobre esse modelo multicamadas,
justifica a escala logarítmica, analisa a fronteira de Pareto, descreve o
algoritmo numérico de otimização (multi-arranque com \texttt{fmincon}) e a
reconstrução da trajetória por interpolação PCHIP, e apresenta um estudo de
sensibilidade ao peso $\alpha$ e ao tamanho da grade de alvos.

O \textbf{Capítulo~\ref{cap:resultados}} (Validação) apresenta a validação
completa por simulação de \textit{backprojection} biestático: configuração do
simulador, cenários, comparação analítico vs.\ simulação, análise de erros e
curvas de projeto.

O \textbf{Capítulo~\ref{cap:conclusoes}} (Conclusões) recapitula as contribuições,
discute as limitações do modelo atual e aponta perspectivas de trabalhos futuros.

O \textbf{Apêndice~\ref{app:math}} reúne as derivações matemáticas de suporte:
conservação de energia nos coeficientes de Fresnel, PSF circular como distribuição
de Bessel, derivação das matrizes ABCD a partir das equações de Maxwell, e
verificação dimensional de $W_z$.

\section{Publicações}
\label{sec:intro_pub}

% TODO: Listar artigos publicados e submetidos originados desta tese.

\textit{[Lista de publicações a ser inserida.]}

%%═══════════════════════════════════════════════════════════════════════════════
\chapter{Fundamentos Teóricos}
\label{cap:fundamentos}
%%═══════════════════════════════════════════════════════════════════════════════

\section{Radar de Abertura Sintética: Princípios e Histórico}
\label{sec:fund_sar_intro}

O SAR é um sistema de imageamento baseado em radar coerente que explora o
movimento da plataforma portadora para sintetizar uma antena virtual de grande
abertura. A técnica, concebida por Carl Wiley em 1951 e demonstrada
experimentalmente por Sherwin et al.\ em 1953~[1], produz imagens de alta
resolução em azimute independentemente do comprimento físico da antena. A
resolução em azimute de um SAR monoestático convencional é $\delta_{az} \approx D/2$,
onde $D$ é o comprimento da antena real, resultado contraintuitivo: a resolução
\emph{melhora} com antenas menores, pois a abertura sintética gerada é inversamente
proporcional a~$D$.

Na configuração biestática, transmissor e receptor encontram-se em plataformas
distintas. Esta separação amplia o espaço de geometrias disponíveis: ângulos de
iluminação e recepção independentes permitem explorar a geometria de espalhamento
do alvo de formas inacessíveis ao radar monoestático~[8,13]. No contexto de SAR
helicoidal, a plataforma — ou as duas plataformas no caso biestático — realiza
trajetórias circulares repetidas em diferentes alturas, gerando cobertura
tridimensional do $k$-espaço da cena.

A tomografia SAR~[16] generaliza o processamento de imagem SAR para múltiplas
passagens, permitindo reconstruir a distribuição tridimensional de espalhadores
em cenas complexas. No regime subsuperficial, a onda deve atravessar a interface
ar-solo, o que introduz a refração e a atenuação por transmissão como fatores
limitantes adicionais do desempenho do sistema~[9,14].

\section{Propagação Eletromagnética em Meios Dielétricos}
\label{sec:fund_em}

\subsection{Equações de Maxwell em Meios Lineares Isotrópicos}
\label{ssec:maxwell}

Considerem-se dois meios dielétricos sem perdas, lineares, isotrópicos e
homogêneos, separados por uma interface plana em $z = 0$. O meio superior
($z > 0$) tem permissividade $\varepsilon_1 = \varepsilon_0$ e permeabilidade
$\mu_1 = \mu_0$ (ar). O meio inferior ($z < 0$) tem
$\varepsilon_2 = \varepsilon_r\varepsilon_0$ e $\mu_2 = \mu_0$ (dielétrico
não-magnético sem perdas). Os índices de refração são $n_1 = 1$ e
$n_2 = \sqrt{\varepsilon_r}$.

As equações de Maxwell no domínio da frequência (notação fasorial com convenção
$e^{\jj\omega t}$) são:

\begin{align}
  \nabla \times \vect{E} &= -\jj\omega\mu\,\vect{H} \label{eq:maxwell_rot_e} \\
  \nabla \times \vect{H} &= +\jj\omega\varepsilon\,\vect{E} \label{eq:maxwell_rot_h} \\
  \nabla \cdot (\varepsilon\vect{E}) &= 0 \label{eq:maxwell_div_e} \\
  \nabla \cdot (\mu\vect{H}) &= 0 \label{eq:maxwell_div_h}
\end{align}

Tomando o rotacional de~\eqref{eq:maxwell_rot_e}, usando~\eqref{eq:maxwell_rot_h}
e a identidade $\nabla\times(\nabla\times\vect{E}) =
\nabla(\nabla\cdot\vect{E}) - \nabla^2\vect{E}$,
juntamente com~\eqref{eq:maxwell_div_e}:

\begin{equation}
  \nabla^2 \vect{E} + k^2 \vect{E} = 0, \qquad k = \frac{\omega n}{c}
  \label{eq:helmholtz}
\end{equation}

\subsection{Ondas Planas e Polarização TM}
\label{ssec:ondas_tm}

Uma onda plana propagando-se no plano $xz$ (plano de incidência) com vetor de
onda $\vect{k} = k(\sin\theta\,\uvec{x} + \cos\theta\,\uvec{z})$ tem a forma
$e^{\jj(\vect{k}\cdot\vect{r} - \omega t)}$. A polarização TM (transverso-magnética)
é definida pela condição em que o campo magnético é transverso ao plano de
incidência: $\vect{H} \parallel \uvec{y}$ e $\vect{E}$ pertence ao plano $xz$.

Para a onda incidente TM, com amplitude $E_0$:
\begin{align}
  \vect{E}^i &= E_0(\cos\theta_i\,\uvec{x} - \sin\theta_i\,\uvec{z})\,
               e^{\jj k_1(x\sin\theta_i + z\cos\theta_i)} \label{eq:E_inc} \\
  \vect{H}^i &= \frac{E_0}{\eta_1}\,\uvec{y}\,
               e^{\jj k_1(x\sin\theta_i + z\cos\theta_i)} \label{eq:H_inc}
\end{align}

onde $\eta_i = \eta_0/n_i$ é a impedância intrínseca do meio $i$. Os campos
refletido e transmitido têm amplitudes $r_{\text{TM}} E_0$ e $t_{\text{TM}} E_0$,
respectivamente.

\subsection{Condições de Contorno e Coeficientes de Fresnel TM}
\label{ssec:fresnel}

As condições de contorno na interface $z=0$ exigem a continuidade das componentes
tangenciais de $\vect{E}$ e $\vect{H}$:

\begin{align}
  E_x^i + E_x^r &= E_x^t \label{eq:bc_ex} \\
  H_y^i + H_y^r &= H_y^t \label{eq:bc_hy}
\end{align}

A condição de fase (princípio de Huygens) impõe que a variação de fase na
direção $x$ seja idêntica para as três ondas em $z=0$:
$k_1\sin\theta_i = k_1\sin\theta_r = k_2\sin\theta_t$, donde $\theta_r = \theta_i$
e a \textbf{lei de Snell}:

\begin{equation}
  \boxed{n_1\sin\theta_i = n_2\sin\theta_t}
  \label{eq:snell}
\end{equation}

Substituindo os campos nas condições~\eqref{eq:bc_ex}--\eqref{eq:bc_hy} e
usando~\eqref{eq:snell}:

\begin{align}
  (1 + r_{\text{TM}})\cos\theta_i &= t_{\text{TM}}\cos\theta_t \label{eq:sys_bc1} \\
  \frac{1}{\eta_1}(1 - r_{\text{TM}}) &= \frac{t_{\text{TM}}}{\eta_2} \label{eq:sys_bc2}
\end{align}

Resolvendo~\eqref{eq:sys_bc1}--\eqref{eq:sys_bc2}, obtêm-se os
\textbf{coeficientes de Fresnel TM}:

\begin{align}
  \boxed{r_{\text{TM}} = \frac{n_2\cos\theta_i - n_1\cos\theta_t}{n_2\cos\theta_i + n_1\cos\theta_t}} \label{eq:r_TM} \\[6pt]
  \boxed{t_{\text{TM}} = \frac{2n_1\cos\theta_i}{n_2\cos\theta_i + n_1\cos\theta_t}} \label{eq:t_TM}
\end{align}

\subsection{Reflectância e Transmitância de Potência}
\label{ssec:pot_fresnel}

O vetor de Poynting médio $\vect{S} = \frac{1}{2}\,\text{Re}(\vect{E}\times\vect{H}^*)$
fornece o fluxo de potência por unidade de área. A componente $z$ para as três
ondas é:

\begin{equation}
  S_z^i = \frac{|E_0|^2}{2\eta_1}\cos\theta_i, \quad
  S_z^r = -\frac{|r_{\text{TM}}|^2|E_0|^2}{2\eta_1}\cos\theta_i, \quad
  S_z^t = \frac{|t_{\text{TM}}|^2|E_0|^2}{2\eta_2}\cos\theta_t
\end{equation}

As grandezas de potência normalizadas são:

\begin{equation}
  R_{\text{TM}} = \frac{|S_z^r|}{S_z^i} = |r_{\text{TM}}|^2
  \label{eq:R_TM}
\end{equation}

\begin{equation}
  T_{\text{TM}} = \frac{S_z^t}{S_z^i} = \frac{n_2\cos\theta_t}{n_1\cos\theta_i}
  \,|t_{\text{TM}}|^2 = 1 - |r_{\text{TM}}|^2
  \label{eq:T_TM}
\end{equation}

A segunda igualdade em~\eqref{eq:T_TM} é a lei de conservação de energia
$R_{\text{TM}} + T_{\text{TM}} = 1$, verificada algebricamente no
Apêndice~\ref{app:energia_fresnel}.

\begin{figure}[H]
\centering
\includegraphics[width=0.85\textwidth]{pv_fig01_fresnel_tm.pdf}
\caption{Reflectância $R_{\text{TM}}$ e transmitância $T_{\text{TM}}$ em função
do ângulo de incidência $\theta_i$ para $n_1=1$ (ar) e $n_2=2$
($\varepsilon_r=4$, solo seco típico). O ângulo de Brewster
$\theta_B \approx 63{,}4^\circ$ é o único zero de $R_{\text{TM}}$.}
\label{fig:fresnel_tm}
\end{figure}

\section{O Ângulo de Brewster}
\label{sec:brewster}

\begin{teorema}[Ângulo de Brewster]
\label{teo:brewster}
Para incidência TM em meios sem perdas e sem atividade magnética
($\mu_1 = \mu_2 = \mu_0$), existe um único ângulo $\theta_B \in (0,\pi/2)$
tal que $r_{\text{TM}}(\theta_B) = 0$, dado por:
\begin{equation}
  \boxed{\theta_B = \arctan\!\left(\frac{n_2}{n_1}\right)}
  \label{eq:brewster}
\end{equation}
\end{teorema}

\begin{proof}
Impondo $r_{\text{TM}} = 0$ em~\eqref{eq:r_TM}: $n_2\cos\theta_B = n_1\cos\theta_t^{(B)}$.
Elevando ao quadrado e usando a lei de Snell $\sin\theta_t^{(B)} = (n_1/n_2)\sin\theta_B$:
\[
  n_2^2 - n_2^2\sin^2\theta_B = n_1^2 - n_1^2\cdot\frac{n_1^2}{n_2^2}\sin^2\theta_B
\]
\[
  \frac{n_2^4-n_1^4}{n_2^2}\sin^2\theta_B = n_2^2 - n_1^2
  \quad\Rightarrow\quad
  \sin^2\theta_B = \frac{n_2^2}{n_1^2+n_2^2}
\]
Portanto $\tan\theta_B = n_2/n_1$.
\end{proof}

\begin{corolario}
\label{cor:brewster_ortho}
No ângulo de Brewster, o raio refratado e o hipotético raio refletido são
perpendiculares: $\theta_B + \theta_t^{(B)} = \pi/2$.
\end{corolario}

\begin{proof}
$\sin\theta_t^{(B)} = (n_1/n_2)\cdot n_2/\sqrt{n_1^2+n_2^2}
= n_1/\sqrt{n_1^2+n_2^2} = \cos\theta_B$,
logo $\theta_t^{(B)} = \pi/2 - \theta_B$.
\end{proof}

O significado físico é imediato: a ortogonalidade implica que os dipolos
induzidos pelos campos refratados oscilam exatamente na direção do hipotético
raio refletido. Como um dipolo não irradia ao longo de seu próprio eixo, a
reflexão se anula por razões intrínsecas ao mecanismo de radiação dipolar~[5].
Para $n_2 = 2$ (solo com $\varepsilon_r = 4$, valor típico de solo seco):
$\theta_B = \arctan(2) \approx 63{,}4^\circ$.

O ângulo de Brewster é o alvo geométrico fundamental do sistema proposto nesta
tese: posicionar o receptor de modo que o raio alvo$\to$Rx incida na interface
ar-solo com $\theta_R \approx \theta_B$ maximiza a transmitância $T_{\text{TM}}$
e, portanto, a potência retornada ao receptor.

\begin{figure}[H]
\centering
\includegraphics[width=0.80\textwidth]{pv_fig08_brewster_geometry.pdf}
\caption{Geometria do ângulo de Brewster. Em $\theta_B$, a ortogonalidade entre
o raio refratado e o hipotético raio refletido impede a radiação na direção
refletida, anulando $r_{\text{TM}}$.}
\label{fig:brewster_geom}
\end{figure}

\section{Traçado de Raios e Princípio de Fermat}
\label{sec:fermat}

\subsection{Formulação Variacional do Princípio de Fermat}
\label{ssec:fermat_var}

O princípio de Fermat estabelece que a trajetória seguida por um raio
eletromagnético entre dois pontos é aquela que torna o caminho óptico
\emph{estacionário}~[5,7].

\begin{definicaom}[Caminho Óptico]
\label{def:caminho_optico}
O caminho óptico entre os pontos $A$ e $B$ ao longo de uma trajetória
$\mathcal{C}$ é:
\begin{equation}
  \mathcal{L}[\mathcal{C}] = \int_{\mathcal{C}} n(\vect{r})\,\dd s
\end{equation}
\end{definicaom}

Para dois meios homogêneos separados pela interface $z=0$, com ponto aéreo
$A = (x_A, y_A, z_A)$ ($z_A > 0$) e ponto subsuperficial $B = (x_B, y_B, z_B)$
($z_B < 0$), o único grau de liberdade geométrico é o ponto de cruzamento
$Q = (u, 0)$ da interface (no plano vertical contendo $A$ e $B$):

\begin{equation}
  \mathcal{L}(u) = n_1\sqrt{u^2 + z_A^2} + n_2\sqrt{(d-u)^2 + z_B^2}
  \label{eq:fermat_L}
\end{equation}

onde $d$ é a separação horizontal projetada entre $A$ e $B$.

\subsection{Condição de Estacionaridade e Lei de Snell}
\label{ssec:fermat_snell}

A condição de estacionaridade $\dd\mathcal{L}/\dd u = 0$ fornece:

\begin{equation}
  \frac{n_1\,u}{\sqrt{u^2+z_A^2}} - \frac{n_2\,(d-u)}{\sqrt{(d-u)^2+z_B^2}} = 0
  \label{eq:fermat_estac}
\end{equation}

que é precisamente $n_1\sin\theta_i = n_2\sin\theta_t$, ou seja, a lei de
Snell~\eqref{eq:snell}. A unicidade do ponto $u^*$ é garantida porque a função
$g(u) = n_1\sin\theta_i(u) - n_2\sin\theta_t(u)$ tem derivada:

\begin{equation}
  \frac{\dd g}{\dd u} = \frac{n_1 z_A^2}{(u^2+z_A^2)^{3/2}}
  + \frac{n_2 z_B^2}{((d-u)^2+z_B^2)^{3/2}} > 0
  \label{eq:g_prime}
\end{equation}

sempre positiva, assegurando que $g$ tem raiz única e que o ponto estacionário
é um mínimo global de $\mathcal{L}$.

\subsection{Método de Newton-Raphson para o Ponto de Refração}
\label{ssec:newton}

A equação $g(u) = 0$ não admite solução de forma fechada para posições 3D
arbitrárias. A iteração de Newton-Raphson:

\begin{equation}
  u^{(k+1)} = u^{(k)} - \frac{g(u^{(k)})}{g'(u^{(k)})}
  \label{eq:NR}
\end{equation}

converge quadraticamente, pois $g'(u) > 0$ garante monotonicidade e existência
de derivada segunda limitada. Na prática, 3--5 iterações com tolerância
$|g| < 10^{-10}$\,m são suficientes para precisão de engenharia.

\subsection{Caminho Geométrico Refratado}
\label{ssec:path_geom}

O comprimento geométrico total do raio Tx$\to$alvo (comprimento físico, não
óptico) é:

\begin{equation}
  \ell = \sqrt{(u^*)^2+z_A^2} + \sqrt{(d-u^*)^2+z_B^2}
  \label{eq:path_geom}
\end{equation}

Este comprimento é usado na equação do radar (perda por espalhamento geométrico
$\propto 1/\ell^2$), enquanto o caminho óptico $\mathcal{L}(u^*)$ determina a
fase do eco recebido e, portanto, o posicionamento do alvo na imagem SAR.

\section{O SAR Biestático Helicoidal}
\label{sec:fund_sar_bist}

\subsection{Geometria da Hélice Cônica}
\label{ssec:helix_geom}

O transmissor descreve uma hélice cônica parametrizada por:

\begin{align}
  x_{Tx}(t) &= \rho(t)\cos\!\left(\frac{2\pi N_v\,t}{T} + \phi_0\right), \nonumber\\
  y_{Tx}(t) &= \rho(t)\sin\!\left(\frac{2\pi N_v\,t}{T} + \phi_0\right), \nonumber\\
  z_{Tx}(t) &= z_{\min} + \frac{(z_{\max}-z_{\min})\,t}{T}
  \label{eq:helix}
\end{align}

onde $\rho(t) = \rho_{\min} + (\rho_{\max}-\rho_{\min})t/T$, $N_v$ é o número
de voltas e $T$ o período total. O receptor voa na mesma hélice, defasado de
um ângulo azimutal fixo $\Delta\phi$ (NorthOffset), de modo que
$\phi_{Rx}(t) = \phi_{Tx}(t) + \Delta\phi$.

Os parâmetros geométricos relevantes para a análise de resolução são:
\begin{itemize}
  \item Ângulo de \textit{look} médio: $\psi_0 = \arctan(\rho_0/z_0)$, onde
        $\rho_0$ e $z_0$ são raio e altura médios.
  \item Comprimento de arco da hélice: $B_{helix} = \sqrt{\Delta\rho^2 + \Delta z^2}$,
        com $\Delta\rho = \rho_{\max}-\rho_{\min}$ e $\Delta z = z_{\max}-z_{\min}$.
  \item Ângulo de inclinação: $\beta = \arctan(\Delta z/\Delta\rho)$.
  \item Apertura tomográfica efetiva: $B_\perp = B_{helix}|\cos(\beta-\psi_0)|$.
\end{itemize}

A \textbf{condição ótima} de inclinação é $\beta = \psi_0$, que maximiza
$B_\perp = B_{helix}$ (apertura tomográfica máxima com dado comprimento de hélice).

\begin{figure}[H]
\centering
\includegraphics[width=0.86\textwidth]{pv_fig10_sistema_3d.pdf}
\caption{Geometria do sistema SAR biestático helicoidal. Tx e Rx descrevem
hélices cônicas ao redor da cena subsuperficial. O NorthOffset $\Delta\phi$ é
o atraso azimutal fixo entre Tx e Rx.}
\label{fig:sistema3d}
\end{figure}

\subsection{Modelo de Sinal e Forma de Onda}
\label{ssec:signal_model}

% TODO: Descrever o modelo de sinal LFM (chirp) transmitido e o sinal
%       recebido após mistura e filtragem por correlação. Incluir a expressão
%       da fase bistática em função das posições Tx, Rx e alvo.
%       Requer derivação matemática nova.

\textit{[Conteúdo a ser desenvolvido.]}

\subsection{Equação do Radar Biestático com Refração}
\label{ssec:radar_eq}

A potência entregue por unidade de área ao alvo pelo transmissor é~[6]:

\begin{equation}
  S_T = \frac{P_t\,G_t}{4\pi\,R_T^2}
  \label{eq:power_density}
\end{equation}

A potência total espalhada pelo alvo é $P_{\text{scat}} = \sigma\,S_T$,
e a densidade de potência no receptor é $S_R = P_{\text{scat}}/(4\pi R_R^2)$.
Com área efetiva de recepção $A_{\text{ef}} = G_r\lambda^2/(4\pi)$, a potência
recebida é:

\begin{equation}
  P_r = \frac{P_t\,G_t\,G_r\,\sigma\,\lambda^2}{(4\pi)^3\,R_T^2\,R_R^2}
  \label{eq:radar_bist}
\end{equation}

Na presença de refração na interface ar-solo, a onda percorre dois tramos em
cada direção: um tramo no ar ($\ell^{(1)}$) e um tramo no dielétrico ($\ell^{(2)}$).
Os coeficientes de transmissão TM da interface modificam a potência efetiva:
$T_1 = T_{\text{TM}}(\theta_T)$ para o caminho Tx$\to$alvo e
$T_2 = T_{\text{TM}}(\theta_R)$ para o caminho alvo$\to$Rx. A equação do radar
biestático com refração é:

\begin{equation}
  \boxed{P_r = \frac{P_t\,G_t\,G_r\,\sigma\,T_1\,T_2\,\lambda^2}{(4\pi)^3\,R_T^2\,R_R^2}}
  \label{eq:radar_refratado}
\end{equation}

onde $R_T$ e $R_R$ são os comprimentos geométricos dos raios refratados dados
por~\eqref{eq:path_geom}, determinados pela lei de Snell via Newton-Raphson.

\subsection{Algoritmo de Backprojection Biestático}
\label{ssec:bp_algo}

% TODO: Descrever o algoritmo de backprojection (BP) biestático com refração:
%       varredura de posições, cálculo do caminho óptico via Snell (Newton-Raphson),
%       acumulação coerente dos ecos. Complexidade computacional O(N_p × N_t).
%       Requer desenvolvimento novo baseado no simulador existente em MATLAB.

\textit{[Conteúdo a ser desenvolvido.]}

\section{Conclusões do Capítulo}
\label{sec:fund_concl}

Este capítulo estabeleceu a base teórica sobre a qual os capítulos subsequentes
são construídos. As equações de Maxwell e as condições de contorno fornecem os
coeficientes de Fresnel TM~\eqref{eq:r_TM}--\eqref{eq:t_TM} e o ângulo de
Brewster~\eqref{eq:brewster}, que é o alvo geométrico do sistema proposto. O
princípio variacional de Fermat demonstra que a lei de Snell é condição
necessária e suficiente para o ponto de refração, com unicidade garantida pela
monotonicidade estrita de $g'(u) > 0$~\eqref{eq:g_prime}. O método de
Newton-Raphson resolve numericamente o problema de refração em 3D com
convergência quadrática.

A geometria da hélice cônica define os parâmetros de apertura tomográfica que
determinam a resolução 3D derivada no Capítulo~\ref{cap:resolucao}. A equação
do radar biestático com refração~\eqref{eq:radar_refratado} é a grandeza a
maximizar no problema de otimização do Capítulo~\ref{cap:otimizacao}.

%%═══════════════════════════════════════════════════════════════════════════════
\chapter{Resolução Espacial Tridimensional do SAR Biestático Helicoidal}
\label{cap:resolucao}
%%═══════════════════════════════════════════════════════════════════════════════

\section{Introdução e Estado da Arte}
\label{sec:res_arte}

% TODO: Revisar a literatura sobre resolução em SAR circular/helicoidal e SAR
%       biestático. Citar Nolan & Fahnestock (2004), Ponce et al. (2014),
%       Knaell & Cardillo (1995), Reigber & Moreira (2000), Góes (2022).
%       Identificar a brecha: nenhum trabalho derivou δ_z e δ_xy para SAR
%       biestático helicoidal com refração na interface ar-solo.

\textit{[Estado da arte a ser desenvolvido.]}

\section{Análise no Espaço de Números de Onda}
\label{sec:kspace}

O fundamento da análise de resolução em SAR é o
\textbf{Teorema da Fatia de Fourier}~[2]:

\begin{teorema}[Munson, O'Brien e Jenkins, 1983]
O processamento SAR em modo \textit{spotlight} é matematicamente equivalente à
tomografia computadorizada: cada posição do radar contribui com uma fatia no
espaço de números de onda da cena. A PSF é a transformada de Fourier inversa
da cobertura acumulada no $k$-espaço.
\end{teorema}

Para o sistema biestático, a fase do eco retornado de um alvo em $\vect{r}_P$
para Tx em $\vect{r}_{Tx}$ e Rx em $\vect{r}_{Rx}$ é~[4]:

\begin{equation}
  \phi = \frac{2\pi f}{c}\!\left(\ell_{T}^{(1)} + n_2\,\ell_{T}^{(2)}
         + \ell_{R}^{(1)} + n_2\,\ell_{R}^{(2)}\right)
  \label{eq:phase_refr}
\end{equation}

onde $\ell^{(1)}$ e $\ell^{(2)}$ denotam os comprimentos de trajeto no ar e
no dielétrico, respectivamente, para cada braço do percurso biestático.

\section{Vetor de Número de Onda Biestático com Refração}
\label{sec:k_bist}

O vetor de número de onda biestático é o gradiente da fase em relação à posição
do alvo:

\begin{equation}
  \vect{k}_{\text{bist}} = \nabla_{\vect{r}_P}\phi
  = -\frac{2\pi f\,n_2}{c}\left(\uvec{\ell}_{T}^{(2)} + \uvec{\ell}_{R}^{(2)}\right)
  \label{eq:k_bist}
\end{equation}

onde $\uvec{\ell}_{T,R}^{(2)}$ são os versores na direção de propagação
\emph{dentro do solo} para os caminhos Tx$\to$alvo e alvo$\to$Rx.
O fator $n_2$ reflete o aumento do número de onda no meio dielétrico.

\section{Resolução Vertical $\delta_z$}
\label{sec:delta_z}

A componente vertical do vetor biestático é:

\begin{equation}
  k_z = \frac{2\pi f\,n_2}{c}\!\left(\cos\theta_T^{(2)} + \cos\theta_R^{(2)}\right)
  \label{eq:kz_bist}
\end{equation}

Para a geometria quasi-monoestática no ponto central da hélice
($\theta_T^{(2)} = \theta_R^{(2)} = \theta_{t,0}$,
$\cos\theta_{t,0} = \sqrt{1-\sin^2\psi_0/n_2^2}$ pela lei de Snell), duas
contribuições independentes estendem o $k$-espaço vertical:

\subsection{Contribuição da Largura de Banda}
\label{ssec:dkz_bw}

Variando $f$ de $f_0$ a $f_0+B$ com posição do sensor fixa:

\begin{equation}
  \Delta k_z^{(B)} = \frac{4\pi n_2 B}{c}\cos\theta_{t,0}
  \label{eq:dkz_bw}
\end{equation}

Este termo existe em qualquer SAR convencional: corresponde à resolução em
alcance refratada. É independente da separação Tx--Rx.

\subsection{Contribuição da Apertura Tomográfica}
\label{ssec:dkz_tom}

Para $f = f_0$ fixo, à medida que o sensor se desloca ao longo da hélice,
$\psi$ varia em torno de $\psi_0$. Usando a derivada
$\dd\cos\theta_{t}/\dd\psi = -\sin\psi/(n_2\cos\theta_t)$ (obtida diferenciando
a lei de Snell refratada) e $|\dd\psi/\dd s| = |\cos(\beta-\psi_0)|/R_0$ para
a hélice cônica:

\begin{equation}
  \Delta k_z^{(\text{tom})} = \frac{4\pi f_0\,B_\perp\,\sin\psi_0\cos\psi_0}{\lambda_0\,R_0\,n_2\cos\theta_{t,0}},
  \qquad B_\perp = B_{helix}|\cos(\beta-\psi_0)|
  \label{eq:dkz_tom}
\end{equation}

\subsection{Fórmula Unificada e Invariância ao NorthOffset}
\label{ssec:delta_z_formula}

A extensão total $\Delta k_z = \Delta k_z^{(B)} + \Delta k_z^{(\text{tom})}$.
Definindo $W_z = c\,\Delta k_z/(4\pi)$, a resolução vertical é:

\begin{equation}
  \boxed{\delta_z = \frac{c}{2\,W_z}}
  \label{eq:delta_z}
\end{equation}

\begin{equation}
  \boxed{W_z = n_2\,B\cos\theta_{t,0}
  + \frac{c\,B_\perp\,\sin\psi_0\cos\psi_0}{\lambda_0\,R_0\,n_2\cos\theta_{t,0}}}
  \label{eq:Wz}
\end{equation}

A resolução vertical~\eqref{eq:delta_z} é \textbf{invariante ao NorthOffset}
$\Delta\phi$, pois a componente $k_z$ do vetor biestático~\eqref{eq:kz_bist}
depende apenas dos ângulos de elevação e não da separação azimutal Tx--Rx.

\section{Resolução Horizontal $\delta_{xy}$}
\label{sec:delta_xy}

Para o $k$-espaço horizontal $(k_x,k_y)$, à medida que Tx percorre $2\pi$ em
azimute e Rx segue defasado de $\Delta\phi$, o vetor biestático varre um anel
de raio:

\begin{equation}
  R_{k} = \frac{2\pi f_0\,\sin\psi_0}{\lambda_0}
  \,\left|\cos\!\left(\frac{\Delta\phi}{2}\right)\right|
  \label{eq:k_ring}
\end{equation}

O fator $|\cos(\Delta\phi/2)|$ expressa a redução da componente horizontal
efetiva pelo NorthOffset biestático: para $\Delta\phi = 0$ (SAR monoestático
equivalente), $R_k$ é máximo; para $\Delta\phi = 180^\circ$ (geometria
anti-podal), $R_k = 0$ e não há resolução horizontal para alvos no eixo.

\begin{figure}[H]
\centering
\includegraphics[width=0.95\textwidth]{fig07_northoffset_kspace.pdf}
\caption{Cobertura do $k$-espaço horizontal para três valores de NorthOffset.
(a) $\Delta\phi=0^\circ$: anel de raio máximo. (b) $\Delta\phi=90^\circ$:
$R_k$ reduzido por $1/\sqrt{2}$. (c) $\Delta\phi=180^\circ$: $R_k=0$, sem
resolução horizontal para alvo no eixo.}
\label{fig:northoffset}
\end{figure}

A distribuição de $k_x = R_k\cos\alpha$ com $\alpha$ uniforme em $[0,2\pi]$
é a distribuição arcseno $p(k_x) = [\pi\sqrt{R_k^2-k_x^2}]^{-1}$, cuja
transformada de Fourier é a função de Bessel $J_0(R_k\delta)$ (demonstração
no Apêndice~\ref{app:bessel}). O critério $-3\,\text{dB}$ ($J_0(u)=1/\sqrt{2}$,
$u_{-3} \approx 1{,}20$) fornece:

\begin{equation}
  \boxed{\delta_{xy}(\Delta\phi) =
  \frac{0{,}60\,\lambda_0}{\pi\,\sin\psi_0\cdot|\cos(\Delta\phi/2)|}}
  \label{eq:delta_xy}
\end{equation}

As expressões~\eqref{eq:delta_z}--\eqref{eq:delta_xy} e~\eqref{eq:Wz},
derivadas nas seções anteriores, constituem o resultado central deste capítulo
e são reproduzidas conjuntamente para facilitar a referência:

\begin{align}
  \delta_z &= \frac{c}{2\,W_z}, \qquad
  W_z = n_2\,B\cos\theta_{t,0}
        + \frac{c\,B_\perp\sin\psi_0\cos\psi_0}{\lambda_0\,R_0\,n_2\cos\theta_{t,0}}
  \label{eq:res_sintese_z} \\[6pt]
  \delta_{xy} &= \frac{0{,}60\,\lambda_0}{\pi\,\sin\psi_0\cdot|\cos(\Delta\phi/2)|}
  \label{eq:res_sintese_xy}
\end{align}

\noindent com $\cos\theta_{t,0} = \sqrt{1-\sin^2\psi_0/n_2^2}$ e
$B_\perp = B_{helix}|\cos(\beta-\psi_0)|$.
A \textbf{condição ótima} de inclinação $\beta = \psi_0$ maximiza $B_\perp = B_{helix}$
e, portanto, minimiza $\delta_z$ para dado comprimento de hélice.

\section{Condição Ótima da Hélice e Correção \textit{Off-Axis}}
\label{sec:helix_opt}

A \textbf{condição ótima} de inclinação $\beta = \psi_0$ maximiza $B_\perp = B_{helix}$,
produzindo a menor $\delta_z$ possível para dado comprimento de hélice. Desvios
$|\beta - \psi_0| > 0$ reduzem $B_\perp$ monotonicamente.

Para um alvo em $x_P \neq 0$ (fora do eixo da hélice), o efeito
\textit{near-range} quebra a simetria cilíndrica. Góes~[3] propõe a correção:

\begin{equation}
  \tilde{\psi}_0 = \arctan\!\left(\frac{\rho_0 - |x_P|}{z_0}\right), \qquad
  \tilde{R}_0 = \sqrt{(\rho_0-|x_P|)^2+z_0^2}
  \label{eq:offaxis}
\end{equation}

nas fórmulas~\eqref{eq:delta_z}--\eqref{eq:delta_xy}, produzindo PSF
assimétrica ($\delta_x \neq \delta_y$). A validação desta correção é apresentada
na Seção~\ref{ssec:val_erros}.

\section{Resultados de Simulação}
\label{sec:res_sim}

% TODO: SIMULAÇÃO — Gerar PSF (função de espalhamento de ponto) via BP
%       biestático e medir δ_z e δ_xy para os seguintes cenários:
%         (a) Δφ = 0°, 45°, 90°, 135°, 180° (variação do NorthOffset)
%         (b) β = ψ_0 (condição ótima) vs. β ≠ ψ_0 (variação de inclinação)
%         (c) n_2 = 1,5 / 2 / 2,5 (variação de permissividade do solo)
%       Plotar: cobertura do k-espaço (k_x,k_y,k_z), cortes da PSF em x, y, z.
%       Comparar analítico vs. simulação em tabela de erros.

\textit{[Resultados de simulação a serem gerados e inseridos.]}

\section{Conclusões do Capítulo}
\label{sec:res_concl}

Este capítulo derivou as equações analíticas de resolução tridimensional do SAR
biestático helicoidal em presença de refração. A resolução vertical
$\delta_z$~\eqref{eq:delta_z} é determinada por duas contribuições aditivas e
independentes — largura de banda e apertura tomográfica — e é invariante ao
NorthOffset. A resolução horizontal $\delta_{xy}$~\eqref{eq:delta_xy} depende
do NorthOffset via o fator $|\cos(\Delta\phi/2)|$ e degrada-se monotonicamente
de $\Delta\phi=0$ a $\Delta\phi = 180^\circ$.

A condição ótima de inclinação $\beta = \psi_0$ maximiza $B_\perp$ e, portanto,
minimiza $\delta_z$ para dado comprimento de hélice. Este resultado, combinado
com a dependência de $\delta_{xy}$ no NorthOffset, define o espaço de projeto
explorado pelo otimizador do Capítulo~\ref{cap:otimizacao}.

%%═══════════════════════════════════════════════════════════════════════════════
\chapter{Transmissão em Meios Estratificados com $N$ Camadas}
\label{cap:multicamadas}
%%═══════════════════════════════════════════════════════════════════════════════

\section{Introdução e Estado da Arte}
\label{sec:multi_arte}

% TODO: Revisar literatura sobre propagação em meios estratificados, formalismo
%       de matrizes de transferência, e aplicações em GPR e subsuperfície.
%       Citar Pozar (2011), Born & Wolf (1999), Daniels (2004), Ishimaru (1991),
%       Banda et al. (2016). Identificar a brecha: aplicação do formalismo ABCD
%       assimétrico (Z_in ≠ Z_out) ao sistema SAR biestático helicoidal.

\textit{[Estado da arte a ser desenvolvido.]}

\section{Geometria do Modelo de $N$ Camadas}
\label{sec:multi_geom}

Considera-se uma pilha de $M$ camadas horizontais. A $i$-ésima camada tem
índice de refração $n_i$ e espessura $d_i$. O meio incidente (ar) tem $n_0=1$.
O alvo encontra-se na camada $L$, à profundidade total
$z_P = -\!\sum_{i=1}^{L-1}d_i - \Delta z_L$, onde $\Delta z_L$ é a
profundidade do alvo dentro da camada $L$. A camada $L$ é tratada como
\textbf{semi-infinita em $z^-$}: não existe interface inferior e toda a energia
transmitida propaga-se para dentro desta camada sem reflexão de retorno.

Esta premissa reflete a situação física de interesse: o alvo encontra-se no
interior de um meio estratificado e não existe uma interface abaixo dele no
domínio de observação. Consequentemente, toda a energia que penetra na camada
$L$ propaga-se em direção ao alvo sem reflexão de retorno pela face inferior,
e a profundidade do alvo dentro dessa camada não é limitada pela espessura
nominal declarada para ela.

\section{O Invariante de Snell Generalizado}
\label{sec:snell_inv}

Em meios estratificados com interfaces horizontais, a componente horizontal
do vetor de onda $k_x = k_i\sin\theta_i = (n_i\omega/c)\sin\theta_i$ é
conservada em cada interface pela condição de fase:

\begin{equation}
  p = n_0\sin\theta_0 = n_1\sin\theta_1 = \cdots = n_M\sin\theta_M
  \label{eq:snell_inv}
\end{equation}

O escalar $p$ é o \emph{parâmetro de raio} (invariante de Snell). O deslocamento
horizontal acumulado pelo raio com parâmetro $p$ até o alvo é:

\begin{equation}
  x(p) = \sum_{i=1}^{L-1} \frac{p\,d_i}{\sqrt{n_i^2-p^2}}
         + \frac{p\,\Delta z_L}{\sqrt{n_L^2-p^2}}
  \label{eq:horiz_p}
\end{equation}

\begin{proposicao}
\label{prop:unicidade_p}
A função $x(p)$ em~\eqref{eq:horiz_p} é estritamente crescente em
$p \in [0, \min_i n_i)$, garantindo a existência e unicidade de $p^*$
tal que $x(p^*) = x_{\text{alvo}}$:
\begin{equation}
  \frac{\dd x}{\dd p} = \sum_{i=1}^{L-1} \frac{n_i^2\,d_i}{(n_i^2-p^2)^{3/2}}
                        + \frac{n_L^2\,\Delta z_L}{(n_L^2-p^2)^{3/2}} > 0
\end{equation}
\end{proposicao}

O parâmetro $p^*$ é determinado por bisseção sobre $f(p)=x(p)-x_{\text{alvo}}=0$,
que converge em $\lceil\log_2((p_{\max}-p_{\min})/\varepsilon)\rceil$ iterações.

\section{Formalismo de Matrizes de Transferência ABCD}
\label{sec:abcd}

A propagação de ondas TM em cada camada é descrita pela equação de Helmholtz:
\begin{equation}
  \frac{\dd^2 H_y}{\dd z^2} + k_{z,i}^2\,H_y = 0, \qquad
  k_{z,i} = \frac{\omega}{c}\sqrt{n_i^2-p^2}
  \label{eq:helmholtz_layer}
\end{equation}

As condições de contorno TM (continuidade de $H_y$ e $(1/n_i^2)\,\partial_z H_y$)
na interface superior e inferior da camada $i$ de espessura $d_i$ são
encapsuladas pela \textbf{matriz de transferência}:

\begin{equation}
  \mathbf{M}_i = \begin{pmatrix}
    \cos(k_{z,i}\,d_i) & \jj\,Z_i\,\sin(k_{z,i}\,d_i) \\[4pt]
    \dfrac{\jj}{Z_i}\,\sin(k_{z,i}\,d_i) & \cos(k_{z,i}\,d_i)
  \end{pmatrix}, \qquad
  Z_i = \frac{n_i}{\cos\theta_i} = \frac{n_i^2}{\sqrt{n_i^2-p^2}}
  \label{eq:ABCD_matrix}
\end{equation}

A derivação de $\mathbf{M}_i$ a partir das equações de Maxwell é apresentada no
Apêndice~\ref{app:abcd_deriv}. Para meios sem perdas, $\det(\mathbf{M}_i) = 1$
(conservação de energia em cada camada). A matriz total das $L-1$ camadas
finitas acima do alvo é:

\begin{equation}
  \begin{pmatrix} A & B \\ C & D \end{pmatrix}
  = \mathbf{M}_1 \cdot \mathbf{M}_2 \cdots \mathbf{M}_{L-1},
  \qquad \det = AD - BC = 1
  \label{eq:ABCD_total}
\end{equation}

\section{Coeficiente de Reflexão e Transmitância Global}
\label{sec:gamma_global}

As impedâncias de entrada e saída da pilha são:

\begin{equation}
  Z_{\text{in}} = \frac{1}{\sqrt{1-p^2}}
  \qquad\text{(ar de entrada)}, \qquad
  Z_{\text{out}} = \frac{n_L^2}{\sqrt{n_L^2-p^2}}
  \qquad\text{(camada do alvo)}
  \label{eq:Zin_Zout}
\end{equation}

Em geral $Z_{\text{in}} \neq Z_{\text{out}}$ pois $n_L > n_0 = 1$ e
$\theta_{\text{out}} < \theta_{\text{in}}$ (lei de Snell). O coeficiente de
reflexão global da pilha, visto do ar de entrada, é:

\begin{equation}
  \boxed{\Gamma_{\text{TM}} = \frac{A Z_{\text{out}} + B
  - Z_{\text{in}}(C Z_{\text{out}} + D)}{%
  A Z_{\text{out}} + B + Z_{\text{in}}(C Z_{\text{out}} + D)}}
  \label{eq:Gamma_TM}
\end{equation}

O coeficiente de transmissão e a transmitância de potência são:

\begin{equation}
  \tau_{\text{TM}} = \frac{2\,Z_{\text{out}}}{A\,Z_{\text{in}} + B
  + C\,Z_{\text{in}}\,Z_{\text{out}} + D\,Z_{\text{out}}}
  \label{eq:tau_TM}
\end{equation}

\begin{equation}
  \boxed{T_{\text{TM}} = |\tau_{\text{TM}}|^2 \cdot \frac{Z_{\text{in}}}{Z_{\text{out}}}
  = 1 - |\Gamma_{\text{TM}}|^2}
  \label{eq:T_TM_multi}
\end{equation}

O fator $Z_{\text{in}}/Z_{\text{out}}$ em~\eqref{eq:T_TM_multi} corrige a
diferença de impedâncias entre os semiespacios de entrada e saída para a
correta contabilização do fluxo de potência~[11].

\begin{lema}
\label{lema:reducao_fresnel}
Para $L=1$ (alvo na primeira camada, sem camadas intermediárias),
a Eq.~\eqref{eq:Gamma_TM} reduz-se exatamente ao coeficiente de Fresnel
TM~\eqref{eq:r_TM}.
\end{lema}

\begin{proof}
$L=1 \Rightarrow \mathbf{M}_{1\cdots 0} = \mathbf{I}$, portanto $[A,B;C,D]=[1,0;0,1]$:
\[
  \Gamma = \frac{Z_{\text{out}} - Z_{\text{in}}}{Z_{\text{out}} + Z_{\text{in}}}
  = \frac{n_2^2/\sqrt{n_2^2-p^2} - 1/\sqrt{1-p^2}}{n_2^2/\sqrt{n_2^2-p^2} + 1/\sqrt{1-p^2}}
  = \frac{n_2/\cos\theta_t - 1/\cos\theta_i}{n_2/\cos\theta_t + 1/\cos\theta_i}
  = \frac{n_2\cos\theta_i - \cos\theta_t}{n_2\cos\theta_i + \cos\theta_t}
  = r_{\text{TM}}
\]
onde foi usada a lei de Snell $p = \sin\theta_i = n_2\sin\theta_t$, de modo que
$\sqrt{n_2^2-p^2} = n_2\cos\theta_t$ e $\sqrt{1-p^2} = \cos\theta_i$.
\end{proof}

\section{Ângulo de Transmissão Ótimo para Pilha Multicamada}
\label{sec:theta_opt}

Para interface única, o máximo de $T_{\text{TM}}$ ocorre em
$\theta_B = \arctan(n_2/n_1)$. Para $M > 1$ camadas, as reflexões múltiplas
nas interfaces internas criam interferências que deslocam o pico de $T_{\text{TM}}$
em relação ao ângulo de Brewster da primeira interface. O ângulo ótimo é
determinado numericamente:

\begin{equation}
  \theta_{\text{opt}} = \arg\max_{\theta_0 \in [0,\pi/2)} T_{\text{TM}}(f,\theta_0;\{n_i,d_i\})
  \label{eq:theta_opt}
\end{equation}

\begin{observacao}[Inacessibilidade do ângulo de Brewster da segunda interface]
Para $n_0=1$, $n_1=\sqrt{2}$, $n_2=2$: o ângulo de Brewster da interface
$n_1$--$n_2$, medido no interior da primeira camada, é
$\theta_B^{(12)} = \arctan(n_2/n_1) = \arctan(\sqrt{2}) \approx 54{,}7^\circ$.
O ângulo correspondente no meio de entrada (ar) seria
$\theta_0 = \arcsin(n_1\sin\theta_B^{(12)}) = \arcsin(\sqrt{2}\cdot\sin 54{,}7^\circ)
= \arcsin(1{,}15) > 1$, valor sem solução real. O ângulo de Brewster da
interface interna é, portanto, inacessível a partir do ar. O máximo de
$T_{\text{TM}}$ para a pilha completa resulta de interferência construtiva entre
as múltiplas reflexões internas e não corresponde a nenhum ângulo de Brewster
clássico de interface isolada.
\end{observacao}

\section{Resultados de Simulação}
\label{sec:multi_sim}

% TODO: SIMULAÇÃO — Calcular T_TM(θ_0) via multislabTMcoef.m para os cenários:
%         (a) Caso de referência: 1 interface (Fresnel puro), n_2 = 2
%         (b) 2 camadas: n_1=1,5, n_2=2,5, d_1 = λ/4 (quarter-wave transformer)
%         (c) 2 camadas: n_1=1,5, n_2=2,5, d_1 = λ/2 (half-wave layer)
%         (d) 3 camadas: perfil de solo estratificado realista
%       Para cada cenário plotar T_TM vs. θ_0 e identificar θ_opt numérico.
%       Verificar conservação de energia R + T = 1 para todos os casos.
%       Script: simulations/run_snell_multislab_pipeline.m

\textit{[Resultados de simulação a serem gerados e inseridos.]}

\section{Conclusões do Capítulo}
\label{sec:multi_concl}

Este capítulo estendeu o modelo eletromagnético de transmissão TM para $N$
camadas estratificadas via formalismo de matrizes de transferência ABCD. O
Lema~\ref{lema:reducao_fresnel} garante que o modelo é uma generalização
consistente do caso de Fresnel, e a Proposição~\ref{prop:unicidade_p} assegura
a unicidade do parâmetro de raio via bisseção. A premissa de última camada
semi-infinita (sem reflexão na face inferior) é adequada para o cenário de
sensoriamento subsuperficial biestático e simplifica a formulação evitando
a necessidade de modelar a geometria abaixo do alvo. A validação numérica da
transmitância $T_{\text{TM}}$ para cenários multicamada é apresentada na
Seção~\ref{sec:val_multi}.

%%═══════════════════════════════════════════════════════════════════════════════
\chapter{Otimização Multi-Objetivo da Trajetória do Receptor}
\label{cap:otimizacao}
%%═══════════════════════════════════════════════════════════════════════════════

\section{Introdução e Estado da Arte}
\label{sec:otim_arte}

% TODO: Revisar literatura sobre otimização de trajetória em SAR biestático,
%       otimização multi-objetivo em sistemas de radar, e uso do ângulo de
%       Brewster para maximizar penetração. Citar Willis (1995), Griffiths (2003),
%       García-Fernández et al. (2019), Flores (2019).
%       Brecha: nenhum trabalho utilizou o ângulo de Brewster como guia
%       geométrico para otimização de trajetória do receptor biestático.

\textit{[Estado da arte a ser desenvolvido.]}

\section{Formulação do Problema de Otimização}
\label{sec:otim_form}

Para uma grade de $N_{tg}$ alvos subsuperficiais, a potência total recebida é,
somando~\eqref{eq:radar_refratado} sobre todos os alvos:

\begin{equation}
  P_r(\vect{r}_{Rx}) = \sum_{i=1}^{N_{tg}}
  \frac{P_t\,G_t^{(i)}\,G_r^{(i)}\,\sigma_i\,T_1^{(i)}\,T_2^{(i)}\,\lambda^2}
       {(4\pi)^3\,R_T^{(i)2}\,R_R^{(i)2}}
  \label{eq:Pr_total}
\end{equation}

Para posição fixa do transmissor $\vect{r}_{Tx}$, o problema de otimização da
posição do receptor $\vect{r}_{Rx}$ no plano horizontal é:

\begin{equation}
  \min_{\vect{r}_{Rx} \in \mathcal{D}}\; J(\vect{r}_{Rx};\alpha) =
  -(1-\alpha)\log_{10}\!\left[P_r(\vect{r}_{Rx})\right]
  + \alpha\,\log_{10}\!\left[\delta_{xy}(\vect{r}_{Rx})\cdot\delta_z(\vect{r}_{Rx})\right]
  \label{eq:cost}
\end{equation}

onde $\mathcal{D} \subset \mathbb{R}^2$ é o domínio de busca e
$\alpha \in [0,1]$ é o parâmetro de compromisso potência--resolução.
O termo $-(1-\alpha)\log_{10}P_r$ penaliza baixa potência recebida (maximiza
$P_r$), enquanto $\alpha\log_{10}(\delta_{xy}\cdot\delta_z)$ penaliza baixa
resolução (minimiza o produto de resoluções).

\section{Justificativa Matemática da Escala Logarítmica}
\label{sec:otim_log}

A escolha da escala logarítmica é motivada por três propriedades:

\begin{enumerate}
  \item \textbf{Comensurabilidade:} $P_r \sim 10^{-15}$\,W e
  $\delta_{xy}\cdot\delta_z \sim 10^{-4}$\,m$^2$ diferem em muitas ordens de
  grandeza. O logaritmo reduz ambos à mesma escala de grandeza.

  \item \textbf{Condicionamento numérico do gradiente:} da equação do radar,
  $-\log_{10}P_r \propto 2\log_{10}R_R - \log_{10}T_2$. O gradiente em relação
  a $\vect{r}_{Rx}$ é proporcional a $\nabla R_R/R_R - \nabla T_2/T_2$,
  que é bem condicionado numericamente para toda a faixa de operação.

  \item \textbf{Interpretação física em dB:} $-10\log_{10}P_r$ é a perda de
  percurso em dB, grandeza de uso corrente em telecomunicações e radar. O
  gradiente de $J$ tem, portanto, interpretação imediata em termos de variação
  de perda de percurso por metro de deslocamento do receptor.
\end{enumerate}

\section{Análise da Fronteira de Pareto}
\label{sec:pareto}

Variando $\alpha$ de 0 a 1, a família de soluções $\vect{r}_{Rx}^*(\alpha)$
parametriza a fronteira de Pareto entre os objetivos $-\log_{10}P_r$ e
$\log_{10}(\delta_{xy}\cdot\delta_z)$.

\begin{proposicao}
Para $\alpha=0$, a solução ótima tende a posicionar o Rx de modo que o ângulo
de incidência $\theta_R$ se aproxime de $\theta_B$ (máximo de $T_2$) e $R_R$
seja mínimo. Para $\alpha=1$, a solução ótima tende ao NorthOffset
$\Delta\phi \to 0$ (máxima resolução horizontal) com mínimo $\psi_0$ (máxima
abertura relativa).
\end{proposicao}

\begin{figure}[H]
\centering
\includegraphics[width=0.82\textwidth]{pv_fig07_pareto.pdf}
\caption{Fronteira de Pareto potência--resolução. Cada ponto corresponde à
posição ótima $\vect{r}_{Rx}^*(\alpha)$ para um valor fixo de $\alpha$.
O joelho da curva em torno de $\alpha \approx 0{,}2$--$0{,}4$ indica a região
de operação recomendada.}
\label{fig:pareto}
\end{figure}

O joelho da curva de Pareto é identificado pelo método da máxima curvatura e
define o intervalo ótimo de compromisso $\alpha^* \in [0{,}2,\, 0{,}4]$.

\section{Pontos de Partida Geometricamente Motivados}
\label{sec:otim_init}

A não-linearidade de $J$ (dependência cúbica em $\vect{r}_{Rx}$ via $R_R^2$ e
$T_2(\theta_R)$) cria múltiplos mínimos locais. Os candidatos iniciais para o
otimizador são geometricamente motivados e fisicamente interpretáveis:

\begin{itemize}
  \item \textbf{Posição oposta:} Rx diametralmente oposto ao Tx em relação ao
  centroide (maximiza $\Delta\phi \approx 180^\circ$, útil para $\alpha \approx 0$
  quando a resolução não é restritiva).

  \item \textbf{Posições perpendiculares:} Rx a $\pm 90^\circ$ do eixo Tx--centroide
  (equilibra cobertura azimutal do $k$-espaço).

  \item \textbf{Posição de Brewster:} Rx à distância horizontal
  $d_B = z_{Rx}\tan\theta_B + |z_P|/\tan\theta_B$ do centroide, de modo que o
  raio alvo$\to$Rx incide na interface com $\theta_R = \theta_B$. O primeiro
  termo é o percurso no ar (ângulo $\theta_B$ em relação à normal); o segundo
  é o percurso no solo, onde o raio refrata para o ângulo complementar
  $\theta_t = \pi/2-\theta_B$ (Corolário~\ref{cor:brewster_ortho}), de modo
  que $\tan\theta_t = 1/\tan\theta_B$. Esta é a posição ótima em potência para
  $\alpha = 0$.
\end{itemize}

\begin{figure}[H]
\centering
\includegraphics[width=0.72\textwidth]{pv_fig06_multistart.pdf}
\caption{Candidatos de arranque \textit{multi-start} em planta $(x,y)$:
posição oposta (1), perpendiculares (2--3) e região de Brewster com
perturbação aleatória (4--5), em torno do centroide dos alvos e dentro da
região de busca de raio $r_s$.}
\label{fig:multistart}
\end{figure}

\section{Algoritmo Numérico de Otimização}
\label{sec:otim_algo}

Para cada posição do transmissor $\vect{r}_{Tx}(t_k)$ ao longo da trajetória
helicoidal, a posição ótima do receptor $\vect{r}_{Rx}^*(t_k)$ que resolve o
problema~\eqref{eq:cost} é obtida por um procedimento de otimização local com
múltiplos arranques (\textit{multi-start}), estruturado nas seguintes etapas:

\begin{enumerate}
  \item \textbf{Decimação da trajetória do Tx.} A hélice completa (milhares
  de amostras, ditadas pela PRF) é subamostrada por um fator de $180$,
  retendo apenas $K$ posições $\{t_k\}$ suficientes para capturar a variação
  lenta da geometria Tx--alvo ao longo do voo. Isso reduz o custo
  computacional de $O(N_{Tx})$ para $O(N_{Tx}/180)$ otimizações completas,
  sem perda de fidelidade — a trajetória completa do Rx é recuperada por
  interpolação (Seção~\ref{sec:pchip}).

  \item \textbf{Pré-cálculo por posição de Tx.} Para cada $t_k$ calculam-se,
  uma única vez, as quantidades que não dependem de $\vect{r}_{Rx}$: os
  ganhos de antena do transmissor para cada alvo $G_t^{(i)}$ e, via o
  invariante de Snell generalizado~\eqref{eq:snell_inv} e o formalismo ABCD
  do Capítulo~\ref{cap:multicamadas}, a distância percorrida $R_T^{(i)}$ e a
  transmitância $T_1^{(i)}$ do enlace Tx$\to$alvo para cada alvo $i$. Essas
  quantidades são reutilizadas em todas as avaliações subsequentes de
  $J$, já que apenas o enlace alvo$\to$Rx varia com $\vect{r}_{Rx}$.

  \item \textbf{Geração de candidatos multi-arranque.} Cinco posições
  iniciais $\vect{x}_0^{(1..5)}$ são geradas segundo a heurística
  geometricamente motivada da Seção~\ref{sec:otim_init} (oposta,
  perpendiculares e Brewster com perturbação aleatória), ilustradas na
  Fig.~\ref{fig:multistart}. Essa diversidade de arranques mitiga o problema
  de múltiplos mínimos locais de $J$ mencionado na Seção~\ref{sec:otim_init}.

  \item \textbf{Otimização local.} A partir de cada candidato $\vect{x}_0^{(i)}$,
  resolve-se
  \[
    \vect{x}^{*(i)} = \arg\min_{\vect{r}_{Rx}\in\mathcal{D}} J(\vect{r}_{Rx};\alpha)
  \]
  com o \textit{solver} \texttt{fmincon} do MATLAB (algoritmo
  \textit{interior-point}), restrito à região quadrada
  $\mathcal{D} = [-r_s,\,r_s]^2$ definida pelo raio de busca $r_s$,
  com tolerância de função $10^{-10}$ e no máximo $200$ iterações.

  \item \textbf{Seleção do melhor arranque.} A posição ótima final é a de
  menor custo entre os cinco resultados locais:
  \[
    \vect{r}_{Rx}^*(t_k) = \vect{x}^{*(i^*)}, \qquad
    i^* = \arg\min_i J\bigl(\vect{x}^{*(i)}\bigr).
  \]
  Essa estratégia \textit{best-of-N} aproxima o mínimo global do problema
  não convexo sem o custo de uma busca exaustiva no domínio $\mathcal{D}$.

  \item \textbf{Recuperação de grandezas físicas.} No ponto ótimo
  $\vect{r}_{Rx}^*(t_k)$, a potência recebida é reavaliada isoladamente com
  $\alpha=0$ na função de custo — desacoplando-a do peso usado na busca — e
  as resoluções $\delta_{xy}$, $\delta_z$ são calculadas pelas fórmulas do
  Capítulo~\ref{cap:resolucao}.
\end{enumerate}

Repetindo esse procedimento para as $K$ posições decimadas do Tx obtém-se a
sequência discreta $\{\vect{r}_{Rx}^*(t_k)\}_{k=1}^{K}$, potência, resolução e
custo associados, que alimentam tanto os resultados da Seção~\ref{sec:otim_sim}
quanto a reconstrução da trajetória contínua a seguir.

\section{Trajetória Completa por Interpolação PCHIP}
\label{sec:pchip}

O otimizador é aplicado a um subconjunto decimado $\{t_k\}_{k=1}^{K}$ de
posições do Tx, produzindo os pontos ótimos $\{\vect{r}_{Rx}^*(t_k)\}$. A
trajetória contínua do Rx é reconstruída por interpolação com polinômios
cúbicos de Hermite preservadores de forma (PCHIP), que:

\begin{itemize}
  \item garantem continuidade de $C^1$ (posição e velocidade contínuas),
  \item preservam monotonia local (sem oscilações de Runge entre pontos),
  \item são realizáveis fisicamente como trajetórias de VANT.
\end{itemize}

A trajetória final é uma hélice deformada que se aproxima da posição de Brewster
do alvo em cada instante, sujeita às restrições de resolução impostas por $\alpha$.

\section{Resultados de Simulação}
\label{sec:otim_sim}

% TODO: SIMULAÇÃO — Complementar com:
%         (b) Posição ótima do Rx (trajetória completa) em planta (x,y)
%         (c) Comparação: trajetória de Brewster puro (α=0) vs. trajetória
%             otimizada (α≈0,3) vs. NorthOffset fixo (referência)
%       Script a usar: simulations/run_snell_multislab_pipeline.m com
%       parâmetros de target_multislab.json.

\subsection{Configuração do Estudo de Sensibilidade}
\label{ssec:otim_sim_config}

Para quantificar o comportamento do otimizador~\eqref{eq:cost} em um cenário
representativo de solo estratificado, executou-se um estudo de sensibilidade
cruzando o peso $\alpha$ com o tamanho da grade de alvos, sobre o pipeline
multicamadas descrito no Capítulo~\ref{cap:multicamadas} (meio de duas
camadas, $n_1=2$, $n_2=3$, $d_1=d_2=1$\,m; portadora em $f=400$\,MHz). O
otimizador foi executado para $\alpha \in \{0;\,0{,}3;\,0{,}6;\,1\}$ e três
grades de alvos de $30$, $60$ e $100$\,m de lado, cada uma com
$5\times5\times5=125$ alvos centrados na segunda camada, totalizando $12$
otimizações completas ao longo de toda a trajetória do transmissor. O objetivo
é caracterizar dois parâmetros de projeto que a análise de Pareto da
Seção~\ref{sec:pareto} não cobre isoladamente: o peso de compromisso $\alpha$
em um meio com múltiplas interfaces, e a extensão espacial da região de
interesse cuja resolução é avaliada.

\subsection{Efeito do Peso $\alpha$ e do Tamanho da Grade de Alvos}
\label{ssec:otim_sim_resultados}

A Tabela~\ref{tab:alpha_grid_resumo} apresenta, para cada combinação de grade
e $\alpha$, os valores médios ao longo da trajetória do Tx da potência
recebida, das resoluções lateral e axial e do produto $\delta_{xy}\cdot\delta_z$
(proxy da área da célula de resolução no plano horizontal).

\begin{table}[H]
\centering
\caption{Métricas médias por cenário do estudo de sensibilidade
$\alpha$--grade (125 alvos, médias ao longo da trajetória do Tx).}
\label{tab:alpha_grid_resumo}
\small
\begin{tabular}{@{}ccrrrr@{}}
\toprule
Grade (m) & $\alpha$ & Potência (dBm) & $\delta_{xy}$ (cm) &
$\delta_z$ (cm) & $\delta_{xy}\!\cdot\!\delta_z$ (cm$^2$) \\
\midrule
30  & 0     & $-24{,}6$ & 1748,7 & 89,8 & $1{,}58\times10^{5}$ \\
30  & 0{,}3 & $-26{,}2$ & 24,4   & 88,8 & 2169 \\
30  & 0{,}6 & $-28{,}2$ & 17,5   & 93,2 & 1628 \\
30  & 1     & $-31{,}0$ & 15,6   & 98,5 & 1540 \\
\midrule
60  & 0     & $-27{,}9$ & 1154,1 & 89,8 & $1{,}04\times10^{5}$ \\
60  & 0{,}3 & $-29{,}2$ & 21,8   & 89,3 & 1943 \\
60  & 0{,}6 & $-31{,}3$ & 17,2   & 93,9 & 1611 \\
60  & 1     & $-33{,}8$ & 15,6   & 98,5 & 1540 \\
\midrule
100 & 0     & $-31{,}0$ & 2196,8 & 89,7 & $1{,}98\times10^{5}$ \\
100 & 0{,}3 & $-32{,}4$ & 25,5   & 89,3 & 2275 \\
100 & 0{,}6 & $-34{,}5$ & 17,1   & 94,3 & 1608 \\
100 & 1     & $-36{,}8$ & 15,6   & 98,5 & 1540 \\
\bottomrule
\end{tabular}
\end{table}

\begin{figure}[H]
\centering
\includegraphics[width=0.82\textwidth]{pv_fig11_tradeoff_multicamada.png}
\caption{Curva de compromisso potência--resolução para o meio de duas
camadas. Cada ponto de uma mesma série (tamanho de grade) corresponde a um
valor de $\alpha$; o eixo horizontal está em escala logarítmica.}
\label{fig:tradeoff_multi}
\end{figure}

\begin{figure}[H]
\centering
\includegraphics[width=0.95\textwidth]{pv_fig12_evolucao_grid30m.png}
\caption{Evolução ao longo da trajetória do Tx para a grade de 30\,m:
potência recebida (topo), produto de resolução $\delta_{xy}\cdot\delta_z$ em
escala logarítmica (meio) e custo $J$ (base), comparando os quatro valores
de $\alpha$.}
\label{fig:evol_grid30}
\end{figure}

\begin{figure}[H]
\centering
\includegraphics[width=0.95\textwidth]{pv_fig13_evolucao_grid100m.png}
\caption{Mesma análise da Fig.~\ref{fig:evol_grid30}, para a grade de
100\,m. O padrão qualitativo em função de $\alpha$ é o mesmo; a diferença
está no nível absoluto de potência.}
\label{fig:evol_grid100}
\end{figure}

O peso $\alpha$ apresenta, no meio multicamadas, o mesmo compromisso
monotônico e fisicamente coerente observado no caso de interface única
(Seção~\ref{sec:pareto}): a correlação entre $\alpha$ e a potência média é
$-1{,}00$ nas três grades (aumentar $\alpha$ \emph{sempre} custa potência), e
a correlação entre $\alpha$ e a resolução (produto $\delta_{xy}\cdot\delta_z$)
é $-0{,}74$ (aumentar $\alpha$ tende a melhorá-la). O compromisso é, no
entanto, fortemente não linear ao longo de $[0,1]$:

\begin{itemize}
  \item De $\alpha=0$ para $\alpha=0{,}3$, o custo em potência é pequeno
        ($1{,}4$--$1{,}6$\,dB, consistente nas três grades), enquanto o
        produto $\delta_{xy}\delta_z$ cai entre $53\times$ e $87\times$ (mais
        de $98\%$ de redução). Esta é a transição de maior retorno: sair de
        $\alpha=0$ — que ignora completamente a resolução, deixando o Rx numa
        geometria com dezenas de metros de erro lateral — custa menos de
        $2$\,dB de potência.
  \item De $\alpha=0{,}3$ para $\alpha=0{,}6$, o custo sobe para
        $\sim\!2$\,dB, e o ganho adicional de resolução cai para
        $17$--$29\%$ — ainda vantajoso, mas com retorno decrescente.
  \item De $\alpha=0{,}6$ para $\alpha=1$, o custo em potência é o maior de
        todos ($2{,}3$--$2{,}9$\,dB) para um ganho de resolução marginal de
        apenas $4$--$5\%$: praticamente todo o benefício de resolução
        alcançável já foi obtido em $\alpha=0{,}6$.
  \item Em $\alpha=1$ a solução converge para a mesma geometria relativa de
        resolução ($\delta_{xy}\approx15{,}6$\,cm, $\delta_z\approx98{,}5$\,cm)
        independentemente da grade — esperado, pois ao ignorar totalmente a
        potência a posição ótima do Rx passa a depender apenas do ângulo
        biestático em relação ao centro da grade, não da sua extensão. Ainda
        assim, a potência física nessa posição varia com o tamanho da grade
        ($-31{,}0 \to -33{,}8 \to -36{,}8$\,dBm), pois alvos mais afastados do
        centro continuam contribuindo menos à potência total mesmo quando a
        posição do Rx não foi escolhida em função deles.
\end{itemize}

Para $\alpha$ fixo, aumentar o tamanho da grade de alvos ($30\to60\to100$\,m)
degrada a potência média recebida com uma penalidade notavelmente uniforme de
$\sim\!3$\,dB por salto de escala de grade, \emph{independentemente do valor
de $\alpha$ usado} — inclusive em $\alpha=1$, quando a posição do Rx é
escolhida ignorando completamente a potência. Isso confirma que o efeito do
tamanho de grade sobre a potência é uma propriedade da distribuição espacial
dos alvos, não um artefato do peso de otimização. Em contraste, o efeito da
grade sobre a resolução é pequeno: para $\alpha$ fixo, o produto
$\delta_{xy}\delta_z$ varia pouco entre grades, pois a resolução é dominada
pela geometria biestática global (posição relativa Tx--Rx--centro da grade),
não pela dispersão espacial dos alvos individuais.

\textbf{Recomendação prática:} para o cenário multicamadas avaliado, a faixa
$\alpha \in [0{,}3,\,0{,}6]$ é o ponto de operação recomendado — captura mais
de $98\%$ do ganho de resolução disponível a um custo de potência de apenas
$1{,}4$--$3{,}6$\,dB frente a $\alpha=0$, e é consistente com o intervalo
$\alpha^*\approx0{,}2$--$0{,}4$ identificado pela fronteira de Pareto de
interface única (Fig.~\ref{fig:pareto}). Empurrar $\alpha$ até $1$ custa mais
$2$--$3$\,dB de potência por um ganho de resolução essencialmente nulo, e por
isso não é recomendado como padrão, exceto quando a resolução lateral for a
única prioridade e houver margem de SNR de sobra. Regiões de interesse mais
extensas (grades maiores) devem ser compensadas com potência de transmissão
adicional ou reprocessadas por sub-regiões menores, já que o tamanho da grade
penaliza a potência de forma consistente sem beneficiar a resolução.

\textbf{Nota sobre a métrica de custo:} o custo $J$ não deve ser comparado
diretamente entre corridas com $\alpha$ distinto, pois sua definição muda com
o peso — os valores brutos de $J$ sugeririam piora monotônica de $\alpha=0$
para $\alpha=1$, o que é apenas um artefato de escala. As comparações de
desempenho desta seção usam sempre as grandezas físicas diretamente
(potência em dBm, $\delta_{xy}$, $\delta_z$), como na Fig.~\ref{fig:tradeoff_multi}.

\section{Conclusões do Capítulo}
\label{sec:otim_concl}

Este capítulo formulou o problema de otimização da trajetória do receptor como
minimização de uma função de custo multi-objetivo logarítmica~\eqref{eq:cost}.
A fronteira de Pareto identifica o intervalo ótimo de compromisso
$\alpha^* \approx 0{,}2$--$0{,}4$, no qual ganhos significativos de potência
são obtidos sem degradação crítica da resolução. A interpolação PCHIP converte
o conjunto de posições ótimas discretas em uma trajetória contínua realizável.
O estudo de sensibilidade da Seção~\ref{sec:otim_sim}, realizado sobre o
pipeline multicamadas, confirma esse intervalo ($\alpha\in[0{,}3,\,0{,}6]$)
para um meio com múltiplas interfaces e mostra, adicionalmente, que o tamanho
da grade de alvos penaliza a potência recebida de forma uniforme
($\sim$3\,dB por salto de escala de grade) sem afetar significativamente a
resolução — um efeito independente do peso $\alpha$ escolhido. Os resultados
numéricos da otimização e a comparação com trajetórias de referência são
apresentados na Seção~\ref{sec:val_otim}.

%%═══════════════════════════════════════════════════════════════════════════════
\chapter{Validação por Simulação de Backprojection}
\label{cap:resultados}
%%═══════════════════════════════════════════════════════════════════════════════

\section{Introdução}
\label{sec:val_intro}

Este capítulo apresenta a validação por simulação de todos os modelos analíticos
derivados nos capítulos anteriores. A estratégia de validação consiste em
comparar previsões analíticas com resultados do simulador de \textit{backprojection}
biestático, que implementa integralmente a física do problema (refração via
Newton-Raphson, fase biestática, acumulação coerente) sem as aproximações usadas
nas derivações analíticas. Erros residuais são interpretados em termos das
aproximações introduzidas explicitamente em cada derivação.

\section{Descrição do Simulador}
\label{sec:val_simulador}

O simulador de \textit{backprojection} biestático integra dois estágios
funcionais sequenciais:

\begin{enumerate}
  \item \textbf{Geração de ecos sintéticos:} para cada par (posição do sensor,
  alvo), o caminho óptico biestático $\mathcal{L}$ é calculado numericamente
  pela resolução iterativa da lei de Snell via Newton-Raphson
  (Seção~\ref{ssec:newton}). A contribuição de fase correspondente
  $e^{-\jj 2\pi f \mathcal{L}/c}$ é acumulada sobre a grade de frequências
  para cada posição da hélice.

  \item \textbf{Retroprojeção tridimensional:} o algoritmo de
  \textit{backprojection} integra coerentemente os ecos gerados no estágio
  anterior sobre todas as posições do sensor, avaliando o campo reconstruído
  em cada ponto da grade tridimensional de imagem. Para o cenário multicamada,
  o caminho óptico em cada tramo é calculado pela bisseção do parâmetro de raio
  $p^*$ conforme a Proposição~\ref{prop:unicidade_p}.
\end{enumerate}

A resolução medida é determinada pela largura a $-3\,\text{dB}$ dos cortes
da função de espalhamento de ponto (PSF) reconstruída nas direções $x$, $y$
e $z$.

\section{Cenários de Validação}
\label{sec:val_cenarios}

Três cenários de validação são definidos:

\begin{description}
  \item[Cenário A — Resolução 3D (interface única):] Alvo pontual em
  $z_P = -5$\,m, solo homogêneo $n_2 = 2$, NorthOffset variável
  $\Delta\phi \in \{90^\circ, 180^\circ\}$, alvo no eixo e fora do eixo
  ($x_P = 20$\,m). Objetivo: validar~\eqref{eq:delta_z} e~\eqref{eq:delta_xy}.

  \item[Cenário B — Otimização da trajetória:] Mesmo alvo do Cenário A, mas
  com trajetória do Rx otimizada para $\alpha \in [0,1]$. Objetivo: validar
  o modelo de potência recebida~\eqref{eq:Pr_total} e confirmar que a trajetória
  ótima se aproxima da posição de Brewster.
  % TODO: SIMULAÇÃO — Executar após implementar o otimizador completo.

  \item[Cenário C — Multicamadas:] Alvo em $z_P = -1{,}5$\,m, pilha de 2
  camadas ($n_1 = 1{,}5$, $d_1 = 0{,}3$\,m; $n_2 = 2{,}5$, semi-infinita).
  Objetivo: validar $\Gamma_{\text{TM}}$~\eqref{eq:Gamma_TM} e
  $T_{\text{TM}}$~\eqref{eq:T_TM_multi}.
  % TODO: SIMULAÇÃO — Executar com parâmetros de target_multislab.json.
\end{description}

\section{Validação das Equações de Resolução 3D}
\label{sec:val_resolucao}

\subsection{Parâmetros da Simulação}
\label{ssec:val_params}

\begin{table}[H]
\centering
\caption{Parâmetros da simulação de validação — Cenário A.}
\label{tab:sim_params}
\begin{tabular}{lrl}
\toprule
Parâmetro & Valor & Descrição \\
\midrule
$f_0$         & 10\,GHz           & Frequência portadora (banda X) \\
$B$           & 50\,MHz           & Largura de banda chirp \\
$n_2$         & 2                 & Índice de refração do solo ($\varepsilon_r=4$) \\
$\rho_0$      & 160\,m            & Raio médio da hélice \\
$z_0$         & 100\,m            & Altura média do sensor \\
$\psi_0$      & 58{,}0$^\circ$    & Ângulo de \textit{look} médio \\
$B_{helix}$   & 47{,}17\,m        & Comprimento total da hélice \\
$\beta$       & 58{,}0$^\circ$    & Ângulo de inclinação ($=\psi_0$: condição ótima) \\
$z_P$         & $-5$\,m           & Profundidade do alvo \\
\bottomrule
\end{tabular}
\end{table}

\subsection{Resultados e Comparação com o Modelo Analítico}
\label{ssec:val_resultados}

\begin{table}[H]
\centering
\caption{Validação do modelo analítico de resolução 3D: previsões vs.\ simulação
de \textit{backprojection} biestático.}
\label{tab:validation}
\begin{tabular}{llrrr}
\toprule
Configuração & Variável & Analítico & Simulação & Erro (\%) \\
\midrule
$\Delta\phi=180^\circ$, no eixo   & $\delta_z$   & 21{,}10\,cm & 21{,}12\,cm & 0{,}1 \\
$\Delta\phi=180^\circ$, no eixo   & $\delta_{xy}$ & 0{,}24\,m   & 0{,}240\,m  & 0{,}0  \\
$\Delta\phi=90^\circ$, no eixo    & $\delta_z$   & 21{,}10\,cm & 21{,}12\,cm & 0{,}1 \\
$\Delta\phi=90^\circ$, no eixo    & $\delta_{xy}$ & 9{,}54\,mm  & 8{,}80\,mm  & 7{,}8  \\
$\Delta\phi=90^\circ$, $x_P=20$\,m & $\delta_z$  & 18{,}72\,cm & 20{,}05\,cm & 7{,}1  \\
$\Delta\phi=90^\circ$, $x_P=20$\,m & $\delta_x$  & 9{,}96\,mm  & 8{,}80\,mm  & 11{,}6 \\
$\Delta\phi=90^\circ$, $x_P=20$\,m & $\delta_y$  & 9{,}96\,mm  & 9{,}10\,mm  & 8{,}6  \\
\bottomrule
\end{tabular}
\end{table}

\noindent O erro percentual é calculado em relação à previsão analítica,
Erro $= |\delta_{\text{analítico}} - \delta_{\text{simulação}}| /
\delta_{\text{analítico}} \times 100\%$, convenção usada de forma consistente
na Tabela~\ref{tab:validation} e na Figura~\ref{fig:validation}.

\begin{figure}[H]
\centering
\includegraphics[width=0.88\textwidth]{fig12_validation.pdf}
\caption{Comparação entre previsões analíticas e resultados de simulação de
\textit{backprojection} biestático para o Cenário A.}
\label{fig:validation}
\end{figure}

\subsection{Análise dos Erros}
\label{ssec:val_erros}

\paragraph{$\delta_z$ no eixo ($0{,}1\%$):}
O erro é compatível com termos de ordem $(B/f_0)^2 \approx 0{,}0025\%$
omitidos na aproximação de banda estreita $B \ll f_0$ usada na derivação
de $\Delta k_z^{(B)}$~\eqref{eq:dkz_bw}. A invariância de $\delta_z$ ao
NorthOffset — mesmo valor para $\Delta\phi = 180^\circ$ e $90^\circ$ — foi
confirmada, validando a independência de $k_z$ com respeito à separação
azimutal.

\paragraph{$\delta_{xy}$ no eixo ($7{,}8\%$):}
A fonte de erro é a aproximação $\psi_0 = \arctan(\rho_0/z_0)$, que ignora
a profundidade do alvo $z_P = -5$\,m. O ângulo efetivo de \textit{look} ao
alvo é $\psi_{\text{ef}} \approx \arctan(\rho_0/|z_0 + n_2|z_P||)$, que
difere de $\psi_0$ em $\sim 1$--$2^\circ$ para estes parâmetros. O erro
cresce linearmente com $|z_P|/z_0$.

\paragraph{$\delta_z$ \textit{off-axis} ($7{,}1\%$):}
A correção de Góes~[3] via~\eqref{eq:offaxis} captura a direção correta
(melhora de $\delta_z$ para $x_P > 0$), mas superestima a magnitude porque
a simulação não pondera os ecos por $1/R^2$. Em um sistema físico real,
o ponderamento natural por $1/R^2$ amplificará o efeito \textit{near-range},
tornando a correção de Góes mais precisa.

\subsection{Curvas de Projeto}
\label{ssec:val_curvas}

\begin{figure}[H]
\centering
\includegraphics[width=0.95\textwidth]{fig10_resolution_curves.pdf}
\caption{Curvas de projeto. (a)~$\delta_z$ vs.\ apertura tomográfica efetiva
$B_\perp$: o termo tomográfico~\eqref{eq:dkz_tom} domina para $B_\perp$
suficientemente grande. (b)~$\delta_{xy}$ vs.\ NorthOffset: degradação severa
em $\Delta\phi \to 180^\circ$. (c)~$\delta_z$ vs.\ ângulo de inclinação $\beta$:
mínimo exato em $\beta = \psi_0$.}
\label{fig:resolution_curves}
\end{figure}

\section{Validação do Modelo de Otimização de Trajetória}
\label{sec:val_otim}

% TODO: SIMULAÇÃO — Após completar o Capítulo de Otimização, executar backprojection com
%       a trajetória otimizada e comparar SNR e resolução obtidos contra
%       trajetórias de referência (NorthOffset fixo, trajetória de Brewster puro).
%       Métrica: ganho em potência recebida (dB) e variação em δ_xy × δ_z (%).

\textit{[Conteúdo a ser desenvolvido após Capítulo~\ref{cap:otimizacao}.]}

\section{Validação do Modelo Multicamadas}
\label{sec:val_multi}

% TODO: SIMULAÇÃO — Executar backprojection biestático com cenário de 2 camadas
%       (n_1=1,5, d_1=0,3 m; n_2=2,5, alvo a 1,5 m de profundidade) e comparar
%       a potência recebida prevista por T_TM(ABCD) com a obtida na simulação.
%       Verificar que o ângulo θ_opt numérico coincide com o pico da curva T_TM.

\textit{[Conteúdo a ser desenvolvido após Capítulo~\ref{cap:multicamadas}.]}

\section{Discussão Geral}
\label{sec:val_discuss}

Os resultados do Cenário A confirmam que o modelo analítico de resolução 3D
derivado no Capítulo~\ref{cap:resolucao} é preciso para alvos no eixo, com
erros inferiores a 0{,}1\% em $\delta_z$. Os erros maiores em $\delta_{xy}$
(7--12\%) para alvos fora do eixo e com $\Delta\phi$ pequeno decorrem de duas
aproximações controladas: (i) o ângulo de \textit{look} efetivo ao alvo
subsuperficial e (ii) a ausência de ponderação por $1/R^2$ na simulação.

A Tabela~\ref{tab:validation} e a Figura~\ref{fig:validation} demonstram que
o modelo analítico é adequado para projeto do sistema (seleção de parâmetros
da hélice, largura de banda e NorthOffset) com precisão suficiente para
engenharia. Refinamentos da correção \textit{off-axis} para alvos profundos
e para ponderação por $1/R^2$ são identificados como trabalho futuro
(Seção~\ref{sec:concl_fut}).

% TODO: Completar esta seção após obter os resultados dos Cenários B e C.

\section{Conclusões do Capítulo}
\label{sec:val_concl}

O simulador de \textit{backprojection} biestático validou as equações analíticas
de resolução 3D com erros aceitáveis para projeto: $< 0{,}1\%$ em $\delta_z$
no eixo e $< 12\%$ em $\delta_{xy}$ fora do eixo. A invariância de $\delta_z$
ao NorthOffset foi confirmada experimentalmente. As validações dos modelos de
otimização (Cenário B) e multicamadas (Cenário C) são identificadas como
requisitos de trabalho futuro, dependentes da conclusão dos capítulos
respectivos.

%%═══════════════════════════════════════════════════════════════════════════════
\chapter{Conclusões e Perspectivas}
\label{cap:conclusoes}
%%═══════════════════════════════════════════════════════════════════════════════

\section{Contribuições Principais}
\label{sec:concl_contrib}

Esta tese apresentou cinco contribuições originais ao campo do SAR biestático
helicoidal para sensoriamento subsuperficial:

\begin{enumerate}

  \item \textbf{Derivação eletromagnética rigorosa dos coeficientes de Fresnel TM
  e do ângulo de Brewster} (Seção~\ref{sec:brewster}): o Teorema~\ref{teo:brewster}
  estabelece que $\theta_B = \arctan(n_2/n_1)$ é o único zero da reflectância TM
  para meios sem perdas, com prova construtiva a partir das equações de Maxwell e
  interpretação física em termos da radiação dipolar.

  \item \textbf{Equações analíticas de resolução 3D para o SAR biestático
  helicoidal com refração} (Seções~\ref{sec:delta_z}--\ref{sec:delta_xy}):
  as fórmulas~\eqref{eq:delta_z} e~\eqref{eq:delta_xy} generalizam resultados
  da literatura para incluir o efeito de refração na interface ar-solo via
  $\cos\theta_{t,0} = \sqrt{1-\sin^2\psi_0/n_2^2}$, e são validadas com
  erros $< 0{,}05\%$ em $\delta_z$ (Tabela~\ref{tab:validation}).

  \item \textbf{Formulação variacional do problema de refração via princípio
  de Fermat} (Seção~\ref{sec:fermat}): a prova de unicidade do ponto de
  refração via monotonicidade estrita de $g'(u) > 0$~\eqref{eq:g_prime}
  fundamenta rigorosamente o algoritmo de Newton-Raphson usado no simulador.

  \item \textbf{Função de custo multi-objetivo logarítmica para otimização
  conjunta de potência e resolução} (Seção~\ref{sec:otim_form}): a
  função~\eqref{eq:cost} resolve o problema de comensurabilidade entre grandezas
  de ordens de magnitude distintas, e a fronteira de Pareto identifica o
  intervalo ótimo $\alpha^* \approx 0{,}2$--$0{,}4$.

  \item \textbf{Extensão do modelo de transmissão TM para $N$ camadas via
  formalismo ABCD com última camada semi-infinita} (Capítulo~\ref{cap:multicamadas}):
  o Lema~\ref{lema:reducao_fresnel} demonstra a consistência com o caso de Fresnel,
  e a Proposição~\ref{prop:unicidade_p} garante a unicidade do parâmetro de raio
  para o traçado de raios em $N$ camadas.

\end{enumerate}

\section{Limitações do Modelo Atual}
\label{sec:concl_lim}

\begin{itemize}

  \item \textbf{Meios sem perdas:} Os modelos de transmissão (Capítulos~\ref{cap:fundamentos}
  e~\ref{cap:multicamadas}) assumem $\varepsilon_r \in \mathbb{R}$. A extensão
  a meios dissipativos (solos úmidos, com $\varepsilon_r \in \mathbb{C}$) tornaria
  $k_{z,i}$ complexo e modificaria o balanço de potência: $R + T < 1$ e
  $\det(\mathbf{M}_i) \neq 1$.

  \item \textbf{Cobertura azimutal completa:} A PSF horizontal tipo $J_0$
  presupõe cobertura azimutal de $2\pi$ completos. Para coberturas parciais, a
  PSF torna-se uma função \textit{sinc} anular, com lóbulos laterais mais elevados.

  \item \textbf{Aproximação de banda estreita:} A resolução vertical~\eqref{eq:delta_z}
  foi derivada com $B \ll f_0$. Para sistemas UWB ($B/f_0 > 0{,}2$), termos de
  segunda ordem em $B/f_0$ devem ser incluídos.

  \item \textbf{Interface plana:} O modelo assume superfície plana e horizontal.
  Superfícies rugosas introduzem espalhamento difuso e fatores de sombreamento
  que modificam os coeficientes de Fresnel efetivos.

  \item \textbf{Validação em ambiente real:} Os resultados deste trabalho são
  baseados exclusivamente em simulação. A validação com dados reais de campo
  é necessária para avaliar o desempenho em condições não-ideais.

\end{itemize}

\section{Perspectivas Futuras}
\label{sec:concl_fut}

\begin{enumerate}

  \item \textbf{Extensão a meios dissipativos:} incorporar
  $\varepsilon_r(f) = \varepsilon'(f) - \jj\varepsilon''(f)$ (modelo de Cole-Cole
  para solos úmidos) e derivar a profundidade de penetração efetiva
  $\delta_{\text{skin}}(f)$ em função da frequência, relevante para a escolha
  ótima de $f_0$.

  \item \textbf{Limites de Cramér-Rao para localização biestática subsuperficial:}
  derivar os limites de variância mínima para a estimação da posição 3D do alvo
  em função de SNR e geometria biestática, fornecendo um limite teórico de
  desempenho independente do algoritmo de processamento.

  \item \textbf{Tratamento de superfícies rugosas:} incorporar o modelo de
  Kirchhoff para interfaces rugosas, que introduz um fator de atenuação de
  coerência $e^{-2k_z^2\sigma_h^2}$ (com $\sigma_h$ a rugosidade rms) nos
  coeficientes de Fresnel efetivos.

  \item \textbf{Validação experimental com dados reais de campo:} implementar
  o sistema em VANTs e realizar campanha de voo sobre alvo subsuperficial
  controlado, comparando a potência recebida e a qualidade da imagem entre
  a trajetória otimizada e a trajetória de referência.

  \item \textbf{Extensão do otimizador para múltiplos alvos a diferentes
  profundidades:} o problema de otimização atual considera um alvo pontual.
  Para cenas com alvos a profundidades distintas, o ângulo de Brewster efetivo
  é diferente para cada alvo, requerendo uma formulação multi-alvo da fronteira
  de Pareto.

\end{enumerate}

%%═══════════════════════════════════════════════════════════════════════════════
\appendix
%%═══════════════════════════════════════════════════════════════════════════════

\chapter{Derivações Matemáticas}
\label{app:math}

\section{Conservação de Energia nos Coeficientes de Fresnel TM}
\label{app:energia_fresnel}

Da condição $S_z^{(1)} = S_z^{(2)}$ (fluxo de potência contínuo na interface):
\[
  \frac{|E_0|^2}{2\eta_1}\cos\theta_i(1-|r_{\text{TM}}|^2)
  = \frac{|t_{\text{TM}}|^2|E_0|^2}{2\eta_2}\cos\theta_t
\]
\[
  T_{\text{TM}} = \frac{n_2\cos\theta_t}{n_1\cos\theta_i}|t_{\text{TM}}|^2
\]

Substituindo~\eqref{eq:r_TM}--\eqref{eq:t_TM}:
\[
  T_{\text{TM}} =
  \frac{n_2\cos\theta_t}{n_1\cos\theta_i}
  \cdot\frac{4n_1^2\cos^2\theta_i}{(n_2\cos\theta_i+n_1\cos\theta_t)^2}
  = \frac{4n_1 n_2\cos\theta_i\cos\theta_t}{(n_2\cos\theta_i+n_1\cos\theta_t)^2}
\]

O balanço $R_{\text{TM}} + T_{\text{TM}}$:
\[
  \frac{(n_2\cos\theta_i-n_1\cos\theta_t)^2 + 4n_1 n_2\cos\theta_i\cos\theta_t}
       {(n_2\cos\theta_i+n_1\cos\theta_t)^2}
  = \frac{(n_2\cos\theta_i+n_1\cos\theta_t)^2}
         {(n_2\cos\theta_i+n_1\cos\theta_t)^2} = 1 \quad\checkmark
\]

\section{PSF Circular como Distribuição de Bessel}
\label{app:bessel}

Seja $\alpha$ uniformemente distribuído em $[0,2\pi]$ e $k_x = R_k\cos\alpha$.
A densidade de probabilidade de $k_x$ é a distribuição arcseno: para
$k_x \in [-R_k, R_k]$,
\[
  p(k_x) = \frac{|\dd\alpha/\dd k_x|}{2\pi/2}
  = \frac{1}{\pi\sqrt{R_k^2-k_x^2}}
\]

A transformada de Fourier é:
\[
  \hat{p}(\delta) = \int_{-R_k}^{R_k}
  \frac{e^{\jj k_x\delta}}{\pi\sqrt{R_k^2-k_x^2}}\dd k_x
  \stackrel{k_x=R_k\cos\varphi}{=}
  \frac{1}{\pi}\int_0^\pi e^{\jj R_k\delta\cos\varphi}\dd\varphi
  = J_0(R_k\delta)
\]

utilizando a representação integral de Bessel
$J_0(z) = \frac{1}{\pi}\int_0^\pi e^{\jj z\cos\varphi}\dd\varphi$.

\section{Derivação das Matrizes ABCD a partir das Equações de Maxwell}
\label{app:abcd_deriv}

Na camada $i$ em polarização TM, a equação de Helmholtz para $H_y$ é:
\begin{equation}
  \frac{\dd^2 H_y}{\dd z^2} + k_{z,i}^2\,H_y = 0, \qquad
  k_{z,i} = \frac{\omega}{c}\sqrt{n_i^2-p^2}
\end{equation}

A solução geral é $H_y(z) = C_i^+e^{\jj k_{z,i}z} + C_i^-e^{-\jj k_{z,i}z}$.
As condições de contorno TM (continuidade de $H_y$ e
$(1/n_i^2)\partial_z H_y = E_x/n_i$) nas faces superior e inferior da camada
são equivalentes à propagação via matriz de transferência. Definindo o vetor de
estado $[H_y, E_x]^T$, a propagação ao longo da camada de espessura $d_i$ é:

\begin{align*}
  H_y(z+d_i) &= H_y(z)\cos(k_{z,i}d_i) + \jj Z_i E_x(z)\sin(k_{z,i}d_i) \\
  E_x(z+d_i) &= \frac{\jj}{Z_i}H_y(z)\sin(k_{z,i}d_i) + E_x(z)\cos(k_{z,i}d_i)
\end{align*}

que é precisamente a matriz $\mathbf{M}_i$ da Eq.~\eqref{eq:ABCD_matrix}.

\section{Verificação Dimensional de $W_z$}
\label{app:dim_wz}

\begin{align*}
  n_2\,B\cos\theta_{t,0}: \quad &
  [\text{adim}]\cdot[\text{Hz}]\cdot[\text{adim}] = \text{Hz} \checkmark \\
  \frac{c\,B_\perp\sin\psi_0\cos\psi_0}{\lambda_0\,R_0\,n_2\cos\theta_{t,0}}: \quad &
  \frac{[\text{m/s}]\cdot[\text{m}]}{[\text{m}]\cdot[\text{m}]\cdot[\text{adim}]}
  = \text{Hz} \checkmark \\
  \delta_z = c/(2W_z): \quad &
  [\text{m/s}]/[\text{Hz}] = \text{m} \checkmark
\end{align*}

\chapter{Lista de Publicações}
\label{app:pub}

% TODO: Inserir lista completa de publicações (artigos em periódico,
%       conferências e relatórios técnicos) originados desta tese.

\textit{[Lista de publicações a ser inserida.]}

%%═══════════════════════════════════════════════════════════════════════════════
\backmatter
%%═══════════════════════════════════════════════════════════════════════════════

\chapter*{Referências}
\addcontentsline{toc}{chapter}{Referências}

% TODO: Expandir para 40–50 referências seguindo padrão IEEE [N].
%       Mínimo recomendado: 55% artigos de periódico, 25% livros, 20% conferências.

\begin{enumerate}

  \item \label{ref:moreira2013}
  A.\ Moreira, P.\ Prats-Iraola, M.\ Younis, G.\ Krieger, I.\ Hajnsek e
  K.\ P.\ Papathanassiou, ``A tutorial on synthetic aperture radar,''
  \textit{IEEE Geoscience and Remote Sensing Magazine}, vol.\ 1, n.\ 1,
  pp.\ 6--43, Mar.\ 2013.

  \item \label{ref:munson1983}
  D.\ C.\ Munson, J.\ D.\ O'Brien e W.\ K.\ Jenkins, ``A tomographic
  formulation of spotlight-mode synthetic aperture radar,''
  \textit{Proceedings of the IEEE}, vol.\ 71, n.\ 8, pp.\ 917--925, Ago.\ 1983.

  \item \label{ref:goes2022}
  J.\ A.\ Góes, \textit{Techniques for High-Resolution 3D Images with
  Synthetic Aperture Radar}, Tese de Doutorado, Universidade Estadual de
  Campinas (UNICAMP), Campinas, SP, Brasil, 2022.

  \item \label{ref:arikan1988}
  O.\ Arikan e D.\ C.\ Munson Jr., ``A tomographic formulation of bistatic
  synthetic aperture radar,'' in \textit{Proc.\ IEEE ICASSP}, pp.\ 1099--1102,
  1988.

  \item \label{ref:bornwolf1999}
  M.\ Born e E.\ Wolf, \textit{Principles of Optics: Electromagnetic Theory
  of Propagation, Interference and Diffraction of Light}, 7ª~ed.,
  Cambridge: Cambridge University Press, 1999.

  \item \label{ref:ulaby1981}
  F.\ T.\ Ulaby, R.\ K.\ Moore e A.\ K.\ Fung, \textit{Microwave Remote
  Sensing: Active and Passive}, vol.\ 1, Reading, MA: Addison-Wesley, 1981.

  \item \label{ref:ishimaru1991}
  A.\ Ishimaru, \textit{Electromagnetic Wave Propagation, Radiation, and
  Scattering}, Englewood Cliffs, NJ: Prentice Hall, 1991.

  \item \label{ref:horne2002}
  A.\ M.\ Horne e G.\ Yates, ``Bistatic synthetic aperture radar,''
  in \textit{Proc.\ IEE Radar Conference}, QinetiQ Malvern, 2002.

  \item \label{ref:garcia2019}
  M.\ García-Fernández, Y.\ López-Fernández e F.\ Las-Heras,
  ``Autonomous airborne 3D SAR imaging system for subsurface sensing,''
  \textit{Remote Sensing}, vol.\ 11, n.\ 20, p.\ 2357, 2019.

  \item \label{ref:banda2016}
  F.\ Banda, J.\ Dall e S.\ Tebaldini,
  ``Single and multipolarimetric P-band SAR tomography of subsurface ice,''
  \textit{IEEE Transactions on Geoscience and Remote Sensing}, vol.\ 54, n.\ 5,
  pp.\ 2832--2845, 2016.

  \item \label{ref:pozar2011}
  D.\ M.\ Pozar, \textit{Microwave Engineering}, 4ª~ed.,
  Hoboken, NJ: Wiley, 2011.

  \item \label{ref:cumming2005}
  I.\ G.\ Cumming e F.\ H.\ Wong, \textit{Digital Processing of Synthetic
  Aperture Radar Data: Algorithms and Implementation},
  Norwood, MA: Artech House, 2005.

  \item \label{ref:willis1995}
  N.\ J.\ Willis, \textit{Bistatic Radar}, 2ª~ed.,
  Norwood, MA: Artech House, 1995.

  \item \label{ref:daniels2004}
  D.\ J.\ Daniels, \textit{Ground Penetrating Radar}, 2ª~ed.,
  London: IET, 2004.

  \item \label{ref:knaell1995}
  K.\ K.\ Knaell e G.\ P.\ Cardillo, ``Radar tomography for the generation of
  three-dimensional images,''
  \textit{IEE Proceedings — Radar, Sonar and Navigation}, vol.\ 142, n.\ 2,
  pp.\ 54--60, Abr.\ 1995.

  \item \label{ref:reigber2000}
  A.\ Reigber e A.\ Moreira, ``First demonstration of airborne SAR tomography
  using multibaseline L-band data,''
  \textit{IEEE Transactions on Geoscience and Remote Sensing}, vol.\ 38, n.\ 5,
  pp.\ 2142--2152, Set.\ 2000.

  \item \label{ref:ponce2014}
  O.\ Ponce, P.\ Prats-Iraola, M.\ Pinheiro, M.\ Rodriguez-Cassola,
  R.\ Scheiber, A.\ Reigber e A.\ Moreira,
  ``Fully polarimetric high-resolution 3-D imaging with circular SAR at L-band,''
  \textit{IEEE Transactions on Geoscience and Remote Sensing}, vol.\ 52, n.\ 6,
  pp.\ 3274--3290, Jun.\ 2014.

  \item \label{ref:zhu2010}
  X.\ X.\ Zhu e R.\ Bamler, ``Very high resolution spaceborne SAR tomography
  in urban environment,''
  \textit{IEEE Transactions on Geoscience and Remote Sensing}, vol.\ 48, n.\ 12,
  pp.\ 4296--4308, Dez.\ 2010.

  \item \label{ref:flores2019}
  A.\ C.\ Flores Rodríguez, \textit{Multipath Influence Analysis on Detection
  Probability of Radars and Optimal Estimators of Velocity in ArcSAR},
  Tese de Doutorado, Universidade Estadual de Campinas (UNICAMP),
  Campinas, SP, Brasil, 2019.

  \item \label{ref:elachi2006}
  C.\ Elachi e J.\ van Zyl, \textit{Introduction to the Physics and Techniques
  of Remote Sensing}, 2ª~ed., Hoboken, NJ: Wiley, 2006.

\end{enumerate}

\end{document}
